\subsection{DIMENSION AND CODIMENSION OF A SUBSPACE OF A VECTOR SPACE}

PROPOSITION 4. Every subspace F of a vector space E is a direct factor of E and

\[
\dim \mathbf {F} + \dim (\mathbf {E} / \mathbf {F}) = \dim \mathbf {E}.\tag{3}
\]

As the quotient vector space E/F is a free module, we know (§ 1, no. 11, Proposition 21) that F is a direct factor of E; relation (3) is then a special case of formula (1) (no. 2).

COROLLARY 1. If E, F, G are vector spaces over a field K, every exact sequence of linear mappings \(0 \rightarrow E \rightarrow F \rightarrow G \rightarrow 0\) splits.

This is another way of expressing Proposition 4 (§ 1, no. 9).

COROLLARY 2. Let \((\mathbf{E}_i)_{0\leqslant i\leqslant n}\) be a finite family of vector spaces over a field K. If there exists an exact sequence of linear mappings

\[
0 \longrightarrow \mathrm{E} _ {0} \stackrel {u _ {0}} {\longrightarrow} \mathrm{E} _ {1} \stackrel {u _ {1}} {\longrightarrow} \mathrm{E} _ {2} \longrightarrow \dots \longrightarrow \mathrm{E} _ {n - 1} \stackrel {u _ {n - 1}} {\longrightarrow} \mathrm{E} _ {n} \stackrel {u _ {n}} {\longrightarrow} 0\tag{4}
\]

the relation

\[
\sum_ {2 k + 1 \leqslant n} \dim \mathrm{E} _ {2 k + 1} = \sum_ {2 k \leqslant n} \dim \mathrm{E} _ {2 k}\tag{5}
\]

holds, or, if all the spaces are finite dimensional,

\[
\sum_ {i = 1} ^ {n} (- 1) ^ {i} \dim \mathrm{E} _ {i} = 0.\tag{6}
\]

Let \(\mathbf{I}_k = \operatorname{Im} u_k = \operatorname{Ker} u_{k+1}\) for \(0 \leqslant k \leqslant n-1\) ; \(\mathbf{I}_{k+1}\) is therefore isomorphic to \(\mathbf{E}_{k+1}/\mathbf{I}_k\) , hence (formula (3)) \(\dim \mathbf{I}_k + \dim \mathbf{I}_{k+1} = \dim \mathbf{E}_{k+1}\) for \(0 \leqslant k \leqslant n-2\) and moreover \(\dim \mathbf{I}_0 = \dim \mathbf{E}_0\) and \(\mathbf{I}_{n-1} = \mathbf{E}_n\) , hence \(\dim \mathbf{I}_{n-1} = \dim \mathbf{E}_n\) . Replacing \(\dim \mathbf{E}_i\) by its expression as a function of the

\[
\sum_ {k = 0} ^ {n - 1}
\]

dim \(I_{k}\) in the two sides of (5), we obtain on each side \(\sum_{k=0}^{n-1}\dim I_{k}\) , whence the corollary.

COROLLARY 3. If M and N are two subspaces of a vector space E, then

\[
\dim (\mathbf {M} + \mathbf {N}) + \dim (\mathbf {M} \cap \mathbf {N}) = \dim \mathbf {M} + \dim \mathbf {N}.\tag{7}
\]

It suffices to apply Corollary 2 to the exact sequence

\[
0 \rightarrow \mathbf {M} \cap \mathbf {N} \rightarrow \mathbf {M} \oplus \mathbf {N} \rightarrow \mathbf {M} + \mathbf {N} \rightarrow 0
\]

(§ 1, no. 8, Proposition 10) taking account of the fact that

\[
\dim (\mathbf {M} \oplus \mathbf {N}) = \dim \mathbf {M} + \dim \mathbf {N}
\]

(no. 2, Proposition 2).

COROLLARY 4. For every subspace \(\mathbf{F}\) of a vector space \(\mathbf{E}\) , \(\dim \mathbf{F} \leqslant \dim \mathbf{E}\) ; if \(\mathbf{E}\) is finite dimensional, the relation \(\dim \mathbf{F} = \dim \mathbf{E}\) is equivalent to \(\mathbf{F} = \mathbf{E}\) .

The first assertion is obvious from (3); further, if dim E is finite, the relation dim F = dim E implies dim(E/F) = 0 by (3) and a vector space of dimension 0 reduces to 0.

COROLLARY 5. If a vector space \(\mathbf{E}\) is the sum of a family \((\mathbf{F}_{\iota})\) of vector subspaces, then

\[
\dim \mathrm{E} \leqslant \sum_ {\iota} \dim \mathrm{F} _ {\iota}.\tag{8}
\]

If further dim E is finite, the two sides of (8) are equal if and only if E is the direct sum of the family \((\mathbf{F}_{\iota})\) .

The inequality (8) follows from (3) and the fact that E is isomorphic to a quotient of \(\bigoplus_{\iota} F_{\iota}\) (§ 1, no. 7, formula (28)). The second assertion is a particular case of § 1, no. 10, Corollary 5 to Proposition 16, for the equality of the two sides of (8) implies that \(\dim F_{\iota} = 0\) except for a finite number of indices.

DEFINITION 2. Given a vector space E, the codimension (with respect to E) of a subspace F of E, denoted by \(\operatorname{codim}_{E} F\) , or simply codim F, is the dimension of E/F (equal to that of any supplement of F in E).

Relation (3) may then be written

\[
\dim \mathrm{F} + \operatorname{codim} \mathrm{F} = \dim \mathrm{E}.\tag{9}
\]

PROPOSITION 5. Let \(\mathbf{F}\) , \(\mathbf{F}'\) be two subspaces of a vector space \(\mathbf{E}\) , such that \(\mathbf{F} \subset \mathbf{F}'\) . Then \(\operatorname{codim}_{\mathbf{E}} \mathbf{F}' \leqslant \operatorname{codim}_{\mathbf{E}} \mathbf{F} \leqslant \dim \mathbf{E}\) . If \(\operatorname{codim}_{\mathbf{E}} \mathbf{F}\) is finite, the relation \(\operatorname{codim}_{\mathbf{E}} \mathbf{F}' = \operatorname{codim}_{\mathbf{E}} \mathbf{F}\) implies \(\mathbf{F} = \mathbf{F}'\) .

The inequality \(\operatorname{codim}_{\mathbf{E}} \mathbf{F} \leqslant \dim \mathbf{E}\) is obvious from (9) and if \(\dim \mathbf{E}\) is finite the relation \(\operatorname{codim}_{E} F = \dim E\) implies \(\dim F = 0\) and hence \(F = \{0\}\) . The remainder of the proposition follows from this, for

\[
\operatorname{codim} _ {\mathbf {E}} \mathbf {F} ^ {\prime} = \operatorname{codim} _ {\mathbf {E} / \mathbf {F}} (\mathbf {F} ^ {\prime} / \mathbf {F}),
\]

since \(\mathbf{E} / \mathbf{F}'\) is canonically isomorphic to \((\mathbf{E} / \mathbf{F}) / (\mathbf{F}' / \mathbf{F})\) (I, § 4, no. 7, Theorem 4).

PROPOSITION 6. If M and N are two subspaces of a vector space E, then

\[
\operatorname{codim} (\mathbf {M} + \mathbf {N}) + \operatorname{codim} (\mathbf {M} \cap \mathbf {N}) = \operatorname{codim} \mathbf {M} + \operatorname{codim} \mathbf {N}.\tag{10}
\]

It suffices to apply Corollary 2 of Proposition 4 to the exact sequence

\[
0 \rightarrow \mathrm{E} / (\mathrm{M} \cap \mathrm{N}) \rightarrow (\mathrm{E} / \mathrm{M}) \oplus (\mathrm{E} / \mathrm{N}) \rightarrow \mathrm{E} / (\mathrm{M} + \mathrm{N}) \rightarrow 0
\]

(§ 1, no. 7, Proposition 10) and use no. 2, Proposition 2.

Note that if \(\mathbf{E}\) is finite dimensional, (10) is a consequence of (7) and (9).

PROPOSITION 7. If \((\mathbf{F}_i)\) is a finite family of subspaces of a vector space \(\mathbf{E}\) , then \(\operatorname{codim}\left(\bigcap_i \mathbf{F}_i\right) \leqslant \sum_i \operatorname{codim} \mathbf{F}_i\) .

If \(\mathbf{F} = \bigcap_{i}\mathbf{F}_{i}\) , \(\mathbf{E} / \mathbf{F}\) is isomorphic to a subspace of the direct sum of the \(\mathbf{E} / \mathbf{F}_i\) (§ 1, no. 7, formula (27)).

Vector subspaces of dimension 1 (resp. dimension 2) of a vector space E are often called lines passing through 0 (resp. planes passing through 0) (or simply lines (resp. planes) if no confusion arises (cf.~§9, no. 3)), by analogy with the language of Classical Geometry; a subspace of E is called a hyperplane passing through 0 (or simply a hyperplane) if it is of codimension 1. Hyperplanes can also be defined as the maximal elements of the set S of vector subspaces of E distinct from E, ordered by inclusion. There is a one-to-one correspondence between the subspaces of E containing a subspace H and the subspaces of E/H (I, §4, no. 7, Theorem 4); if E is of dimension ≥1, S is non-empty and to say that H is maximal in S means that E/H contains no subspace distinct from \{0\} and E/H, which implies that E/H is generated by any of its elements ≠0, in other words it is of dimension 1.

In a vector space of finite dimension \(n \geqslant 1\) , the hyperplanes are the subspaces of dimension n - 1 by formula (3).

PROPOSITION 8. In a vector space E over a field K, every vector subspace F is the intersection of the hyperplanes which contain it.

It suffices to show that for all \(x \notin F\) there exists a hyperplane H containing F and not containing x. By hypothesis \(F \cap Kx = \{0\}\) and hence the sum M of F and Kx is direct. Let N be supplementary to M in E; E is then the direct sum of \(H = F + N\) and Kx and H is therefore a hyperplane with the desired property.

Remark. Most of the properties proved in this no. for subspaces of a vector space do not hold for submodules of a free module whose dimension (§ 2, no. 7, Remark 2) is defined. *For example, an ideal of a commutative ring does not necessarily admit a basis, for there are integral domains A in which certain ideals are not principal (VII, § 1, no. 1) and any two elements of such a ring are linearly dependent (§ 1, no. 11, Remark 1).* A submodule of a free A-module E may be free, distinct from E and have the same dimension as E, as is shown by principal ideals in an integral domain A; the same example proves moreover that a free submodule of a free A-module does not necessarily admit a supplement.

\subsection{RANK OF A LINEAR MAPPING}

DEFINITION 3. Let E, F be two vector spaces over a field K. For every linear mapping u of E into F, the dimension of the subspace u(E) of F is called the rank of u and denoted by rg(u).

If \(\mathbf{N} = \operatorname{Ker}(u)\) , \(\mathbf{E} / \mathbf{N}\) is isomorphic to \(u(\mathbf{E})\) , whence the relation

\[
\operatorname{rg} (u) = \operatorname{codim} _ {\mathrm{E}} (\operatorname{Ker} (u))\tag{11}
\]

and therefore

\[
\operatorname{rg} (u) + \dim (\operatorname{Ker} (u) = \dim E.\tag{12}
\]

Moreover, by formula (3)

\[
\operatorname{rg} (u) + \dim (\operatorname{Coker} (u)) = \dim F.\tag{13}
\]

PROPOSITION 9. Let E, F be two vector spaces over a field K and u:E→F a linear mapping.

\begin{enumerate}
\def\labelenumi{(\roman{enumi})}
\item
  \(\operatorname{rg}(u) \leqslant \inf(\dim E, \dim F)\) .
\item
  Suppose that \(\mathbf{E}\) is finite-dimensional; in order that \(\operatorname{rg}(u) = \dim \mathbf{E}\) , it is necessary and sufficient that \(u\) be injective.
\item
  Suppose that \(\mathbf{F}\) is finite-dimensional; in order that \(\operatorname{rg}(u) = \dim \mathbf{F}\) , it is necessary and sufficient that \(u\) be surjective.
\end{enumerate}

This follows immediately from relations (12) and (13).

COROLLARY. Let E be a vector space of finite dimension n and u an endomorphism of E. The following properties are equivalent:

\begin{enumerate}
\def\labelenumi{(\alph{enumi})}
\item
  \(u\) is bijective;
\item
  \(u\) is injective;
\item
  \(u\) is surjective;
\item
  \(u\) is right invertible;
\item
  \(u\) is left invertible;
\item
  \(u\) is of rank \(n\) .
\end{enumerate}

If E is an infinite-dimensional vector space, there are injective (resp. surjective) endomorphisms of E which are not bijective (Exercise 9).

Let K, \(K'\) be two fields, \(\sigma:K\to K'\) an isomorphism of K onto \(K'\) , E a vector K-space, \(E'\) a vector \(K'\) -space and \(u:K\to K'\) a semi-linear mapping relative to \(\sigma\) (§1, no. 13); the dimension of the subspace \(u(\mathbf{E})\) of \(E'\) is also called the rank of u. It is also the rank of u considered as a linear mapping of E into \(\sigma_{*}(\mathbf{E}')\) , for every basis of \(u(\mathbf{E})\) is also a basis of \(\sigma_{*}(u(\mathbf{E}))\) .

\subsection{DUAL OF A VECTOR SPACE}

THEOREM 4. The dimension of the dual \(\mathbf{E}^*\) of a vector space \(\mathbf{E}\) is at least equal to the dimension of \(\mathbf{E}\) . For \(\mathbf{E}^*\) to be finite-dimensional, it is necessary and sufficient that \(\mathbf{E}\) be so, and then \(\dim \mathbf{E}^* = \dim \mathbf{E}\) .

If K is the field of scalars of E, E is isomorphic to a space \(\mathbf{K}_{s}^{\mathrm{(I)}}\) and therefore \(E^{*}\) is isomorphic to \(K_{d}^{I}\) (§2, no. 6, Proposition 10). As \(\mathbf{K}_{s}^{\mathrm{(I)}}\) is a subspace of \(K_{d}^{I}\) , \(\dim E = \text{Card}(I) \leqslant \dim E^{*}\) (no. 3, Corollary 4 to Proposition 4); further, if I is finite, \(\mathbf{K}_{d}^{I} = \mathbf{K}_{d}^{\mathrm{(I)}}\) (cf.~Exercise 3(d)).

COROLLARY. For a vector space \(\mathbf{E}\) , the relations \(\mathbf{E} = \{0\}\) and \(\mathbf{E}^{*} = \{0\}\) are equivalent.

THEOREM 5. Given two exact sequences of vector spaces (over the same field K) and linear mappings

\[
0 \rightarrow \mathrm{E} ^ {\prime} \rightarrow \mathrm{E} \rightarrow \mathrm{E} ^ {\prime \prime} \rightarrow 0
\]

\[
0 \to \mathbf {F} ^ {\prime} \to \mathbf {F} \to \mathbf {F} ^ {\prime \prime} \to 0
\]

and two vector spaces G, H over K, the corresponding sequences

\[
0 \rightarrow \operatorname{Hom} (\mathrm{E} ^ {\prime \prime}, \mathrm{G}) \rightarrow \operatorname{Hom} (\mathrm{E}, \mathrm{G}) \rightarrow \operatorname{Hom} (\mathrm{E} ^ {\prime}, \mathrm{G}) \rightarrow 0
\]

\[
0 \rightarrow \operatorname{Hom} (\mathrm{H}, \mathrm{F} ^ {\prime}) \rightarrow \operatorname{Hom} (\mathrm{H}, \mathrm{F}) \rightarrow \operatorname{Hom} (\mathrm{H}, \mathrm{F} ^ {\prime \prime}) \rightarrow 0
\]

are exact and split.

This follows from the fact that every vector subspace is a direct factor (no. 3, Proposition 4) and from § 2, no. 1, Propositions 1 and 2.

COROLLARY. For every exact sequence

\[
0 \longrightarrow \mathrm{E} ^ {\prime} \stackrel {u} {\longrightarrow} \mathrm{E} \stackrel {v} {\longrightarrow} \mathrm{E} ^ {\prime \prime} \longrightarrow 0
\]

of vector spaces over the same field K and linear mappings the sequence

\[
0 \longrightarrow \mathrm{E} ^ {\prime \prime *} \stackrel {t _ {v}} {\longrightarrow} \mathrm{E} ^ {*} \stackrel {t _ {u}} {\longrightarrow} \mathrm{E} ^ {\prime *} \longrightarrow 0
\]

is exact and splits.

It follows in particular that for every vector subspace M of E the canonical homomorphism \(E^{*}/M^{\prime} \rightarrow M^{*}\) , where \(M^{\prime}\) is the subspace of \(E^{*}\) orthogonal to M ( \(\S 2\) , no. 4), is bijective.

THEOREM 6. For every vector space E over a field K, the canonical mapping \(c_{\mathbf{E}}: \mathbf{E} \to \mathbf{E}^{**}\) (§2, no. 7) is injective; for it to be bijective, it is necessary and sufficient that E be finite-dimensional.

The first assertion and the fact that if E is finite-dimensional \(c_{E}\) is bijective are special cases of §2, no. 7, Proposition 14. Suppose that E is infinite-dimensional, so that we may assume that \(\mathbf{E} = \mathbf{K}_{s}^{(L)}\) , where L is an infinite set and therefore \(E^{*} = K_{d}^{L}\) . Let \((e_{\lambda})_{\lambda \in L}\) be the canonical basis of E and let \((e_{\lambda}^{*})_{\lambda \in L}\) be the corresponding family of coordinate forms in \(E^{*}\) (§2, no. 6); the vector subspace of \(E^{*}\) generated by the \(e_{\lambda}^{*}\) is just the direct sum \(F' = K_{d}^{(L)}\) and the hypothesis that L is infinite implies that \(F' \neq E^{*}\) . Then there exists a hyperplane \(H'\) of \(E^{*}\) containing \(F'\) (no. 3, Proposition 8) and, as \(E^{*}/H'\) is non-zero, its dual is also non-zero (Corollary to Theorem 4), which is identified with the orthogonal \(H''\) of \(H'\) in \(E^{**}\) (§2, no. 6, Corollary to Proposition 9). But \(H'' \cap c_{E}(E)\) is contained in the image under \(c_{E}\) of the orthogonal of \(F'\) in E, which is by definition 0; hence \(c_{E}(E) = E^{**}\) is impossible.

E will usually be identified with the subspace of \(E^{**}\) the image of \(c_{E}\) .

Let \(\mathbf{E}, \mathbf{F}\) be two vector spaces over a field \(\mathbf{K}\) and \(u: \mathbf{E} \to \mathbf{F}\) a linear mapping. We shall define canonical isomorphisms:

\begin{enumerate}
\def\labelenumi{(\arabic{enumi})}
\item
  Of the dual of \(\operatorname{Im}(u) = u(\mathbf{E})\) onto \(\operatorname{Im}(^t u) = {}^t u(\mathbf{F}^*)\) .
\item
  Of the dual of \(\operatorname{Ker}(u) = \overline{u}^1(0)\) onto \(\operatorname{Coker}(^t u) = \mathrm{E}^* / {}^t u(\mathrm{F}^*)\) .
\item
  Of the dual of \(\operatorname{Coker}(u) = \mathbf{F} / u(\mathbf{E})\) onto \(\operatorname{Ker}(^t u) = {}^t u^{-1}(0)\) .
\end{enumerate}

We write \(\mathbf{I} = \operatorname{Im}(u),\mathbf{N} = \operatorname{Ker}(u),\mathbf{C} = \operatorname{Coker}(u)\) ; from the exact sequences

\[
0 \longrightarrow \mathrm{N} \longrightarrow \mathrm{E} \stackrel {p} {\longrightarrow} \mathrm{I} \longrightarrow 0, \quad 0 \longrightarrow \mathrm{I} \stackrel {j} {\longrightarrow} \mathrm{F} \longrightarrow \mathrm{C} \longrightarrow 0 \tag {14}
\]

we derive, by transposition (Corollary to Theorem 5), the exact sequences

\[
\begin{array}{l} 0 \longrightarrow \mathrm{I} ^ {*} \stackrel {{t _ {p}}} {{\longrightarrow}} \mathrm{E} ^ {*} \longrightarrow \mathrm{N} ^ {*} \longrightarrow 0, \\ 0 \longrightarrow \mathrm{C} ^ {*} \longrightarrow \mathrm{F} ^ {*} \stackrel {{t _ {j}}} {{\longrightarrow}} \mathrm{I} ^ {*} \longrightarrow 0. \end{array}\tag{15}
\]

Moreover, as \(u = j \circ p\) , \(^{t}u = ^{t}p \circ ^{t}j\) ; the exact sequences (15) thus define canonical isomorphisms of \(C^*\) onto \(\operatorname{Ker}(^{t}u)\) , of \(I^*\) onto \(\operatorname{Im}(^{t}u)\) and of \(N^*\) onto \(\operatorname{Coker}(^{t}u)\) , since \(^{t}p\) is injective and \(^{t}j\) surjective. To be precise, let \(y \in \operatorname{Im}(u)\) , \(z \in \operatorname{Ker}(u)\) , \(t \in \operatorname{Coker}(u)\) , \(y' \in \operatorname{Im}(^{t}u)\) , \(z' \in \operatorname{Coker}(^{t}u)\) , \(t' \in \operatorname{Ker}(^{t}u)\) ; when \(y'\) , \(z'\) , \(t'\) are canonically identified with linear forms on \(\operatorname{Im}(u)\) , \(\operatorname{Ker}(u)\) and \(\operatorname{Coker}(u)\) respectively, then

\[
\langle y, y ^ {\prime} \rangle = \langle x, y ^ {\prime} \rangle \quad \text { for   all } x \in \mathrm{E} \text { such   that } u (x) = y;\tag{16}
\]

\begin{enumerate}
\def\labelenumi{(\arabic{enumi})}
\setcounter{enumi}{16}
\item
  \(\langle z, z' \rangle = \langle z, x^* \rangle\) for all \(x^* \in \mathbf{E}^*\) whose class mod. \(^t u(\mathbf{F})\) is equal to \(z'\) ;
\item
  \(\langle t, t' \rangle = \langle s, t' \rangle\) for all \(s \in \mathbf{F}\) whose class mod. \(u(\mathbf{E})\) is equal to \(t\) .
\end{enumerate}

In particular we derive from these results:

PROPOSITION 10. Let E, F be two vector spaces over the same field K and \(u: \mathbf{E} \to \mathbf{F}\) a linear mapping.

\begin{enumerate}
\def\labelenumi{(\roman{enumi})}
\item
  For \(u\) to be injective (resp. surjective), it is necessary and sufficient that \(^{t}u\) be surjective (resp. injective).
\item
  \(\operatorname{rg}(u) \leqslant \operatorname{rg}(^t u)\) and \(\operatorname{rg}(u) = \operatorname{rg}(^t u)\) if \(\operatorname{rg}(u)\) is finite.
\end{enumerate}

The second assertion follows from the above and Theorem 4.

THEOREM 7. Let E be a vector space over a field K, F a subspace of E and F' the orthogonal of F in E*.

\begin{enumerate}
\def\labelenumi{(\roman{enumi})}
\item
  \(\dim \mathbf{F}' \geqslant \operatorname{codim}_{\mathbb{E}} \mathbf{F}\) ; for \(\dim \mathbf{F}'\) to be finite, it is necessary and sufficient that \(\operatorname{codim}_{\mathbb{E}} \mathbf{F}\) be finite, and then \(\dim \mathbf{F}' = \operatorname{codim}_{\mathbb{E}} \mathbf{F}\) .
\item
  The orthogonal of \(\mathbf{F}'\) in \(\mathbf{E}\) is equal to \(\mathbf{F}\) .
\item
  Every finite-dimensional subspace \(\mathbf{G}'\) of \(\mathbf{E}^*\) is the orthogonal of some subspace of \(\mathbf{E}\) , necessarily equal to the orthogonal of \(\mathbf{G}'\) in \(\mathbf{E}\) and of finite codimension.
\item
  We know that \(\mathbf{F}'\) is isomorphic to the dual \((\mathbf{E} / \mathbf{F})^*\) (§ 2, no. 6, Corollary to Proposition 9) and hence the assertion follows from Theorem 4, since \(\dim (\mathbf{E} / \mathbf{F}) = \operatorname{codim}_{\mathbf{E}}\mathbf{F}\) by definition.
\item
  Let \(F_{1}\) be the orthogonal of \(F'\) in E; clearly \(F \subset F_{1}\) and the orthogonal \(F_{1}'\) of \(F_{1}\) is equal to \(F'\) (§2, no. 4); the canonical linear mapping \((\mathrm{E}/\mathrm{F}_{1})^{*} \to (\mathrm{E}/\mathrm{F})^{*}\) , the transpose of \(E/F \to E/F_{1}\) is therefore bijective (§2, no. 6, Corollary to Proposition 10); it then follows from Proposition 10 that the canonical mapping \(E/F \to E/F_{1}\) is bijective, which implies \(F_{1} = F\) .
\item
  Let \(\mathbf{G}'\) be a subspace of \(\mathbf{E}^*\) of finite dimension \(p\) and let \(\mathbf{F}\) be its orthogonal in \(\mathbf{E}\) ; then \(\operatorname{codim}_{\mathbf{E}} \mathbf{F} \leqslant \dim \mathbf{G}'\) . For, if \((a_{i}^{*})_{1 \leqslant i \leqslant p}\) is a basis of \(\mathbf{G}'\) , \(\mathbf{F}\) is the kernel of the linear mapping \(x \mapsto (\langle x, a_{i}^{*} \rangle)\) from \(\mathbf{E}\) to \(\mathbf{K}_{s}^{p}\) whose rank is at most \(p\) (no. 4, Proposition 9), whence the conclusion (no. 4). Then let \(\mathbf{F}'\) be the orthogonal of \(\mathbf{F}\) in \(\mathbf{E}^*\) ; it follows from (i) that \(\dim \mathbf{F}' \leqslant \dim \mathbf{G}'\) ; but on the other hand obviously \(\mathbf{G}' \subset \mathbf{F}'\) , whence \(\mathbf{F}' = \mathbf{G}'\) ( \(\S 2\) , no. 3, Corollary 4 to Proposition 4).
\end{enumerate}

Remark. An infinite-dimensional subspace \(G'\) of \(E^{*}\) is not necessarily the orthogonal of a subspace of E, in other words, if F is the orthogonal of \(G'\) in E, the orthogonal \(F'\) of F in \(E^{*}\) may be distinct from \(G'\) (Exercise 20(b)).†

COROLLARY 1. Let \((x_{i}^{*})_{1\leqslant i\leqslant p}\) be a finite sequence of linear forms on E and let F be the subspace of E consisting of the \(x\) such that

\[
\langle x, x _ {i} ^ {*} \rangle = 0 \quad for 1 \leqslant i \leqslant p.
\]

Then \(\operatorname{codim}_{E} F\) is equal to the rank of the set of the \(x_{i}^{*}\) and every linear form on E which is zero on F is a linear combination of the \(x_{i}^{*}\) . Then \(\operatorname{codim}_{E} F \leqslant p\) and in order that \(\operatorname{codim}_{E} F = p\) , it is necessary and sufficient that the \(x_{i}^{*}\) be linearly independent.

The set \(G'\) of linear combinations of the \(x_{i}^{*}\) is a subspace of \(E^{*}\) and F is the orthogonal of \(G'\) in E, hence \(\operatorname{codim}_{E} F = \dim G'\) by Theorem 7; further \(\dim G' \leqslant p\) and the relation \(\dim G' = p\) means that \((x_{i}^{*})\) is a free system (no. 2, Proposition 1); hence the corollary.

COROLLARY 2. (i) Let \((x_{i}^{*})_{1\leqslant i\leqslant p}\) be a finite sequence of linear forms on E. For \((x_{i}^{*})\) to be a free system, it is necessary and sufficient that there exist a sequence \((x_{i})_{1\leqslant i\leqslant p}\) of elements of E such that \(\langle x_i,x_j^*\rangle = \delta_{ij}\) (Kronecker index).

\begin{enumerate}
\def\labelenumi{(\roman{enumi})}
\setcounter{enumi}{1}
\tightlist
\item
  Let \((x_{i})_{1\leqslant i\leqslant p}\) be a finite sequence of elements of E. For \((x_{i})\) to be a free system, it is necessary and sufficient that there exist a sequence \((x_{i}^{*})_{1\leqslant i\leqslant p}\) of linear forms on E such that \(\langle x_i,x_j^*\rangle = \delta_{ij}\) .
\end{enumerate}

Clearly (ii) follows from (i) by considering E as identified with a subspace of \(\mathbf{E}^{**}\) by means of \(c_{\mathbf{E}}\) (Theorem 6). Let \(\mathbf{G}'\) be the subspace of \(\mathbf{E}^*\) generated by the \(x_{i}^{*}\) and \(\mathbf{F}\) its orthogonal in \(\mathbf{E}\) ; \(\mathbf{E} / \mathbf{F}\) and \(\mathbf{G}'\) can each be canonically identified with the dual of the other; if the family \((x_{i}^{*})\) is free, there is in \(\mathbf{E} / \mathbf{F}\) a basis \((\dot{x}_i)\) the dual of \((x_{i}^{*})\) and every representative system \((x_{i})\) of the classes \(\dot{x}_{i}\) has the required properties. Conversely, the existence of the system \((x_{i})\) such that \(\langle x_{i}, x_{j}^{*} \rangle = \delta_{ij}\) implies that for all \(i\) the subspace of \(\mathbf{E}^*\) orthogonal to \(\mathbf{K}\) . \(x_{i}\) contains the \(x_{j}^{*}\) of index \(j \neq i\) but does not contain \(x_{i}^{*}\) , hence the system \((x_{i}^{*})_{1 \leqslant i \leqslant p}\) is free.

COROLLARY 3. Let S be a set and V a vector subspace of the right vector K-space \(K_{d}^{S}\) of mappings of S into K. In order that \(\dim V \geqslant p\) (where p is an integer), it is necessary and sufficient that there exist p elements \(s_{i}\) of S and p elements \(f_{i}\) of V ( \(1 \leqslant i \leqslant p\) ) such that \(f_{i}(s_{j}) = \delta_{ij}\) .

The space \(K_{d}^{S}\) is canonically identified with the dual of \(E = K_{d}^{(S)}\) and \(f(s) = \langle e_{s}, f \rangle\) for \(s \in S\) and \(f \in K_{s}^{S}\) , \((e_{s})_{s \in S}\) being the canonical basis of E. Corollary 2 then shows that the condition is sufficient. Conversely, suppose that \(\dim V \geqslant p\) , so that there exists a subspace \(G'\) of V of dimension p; let F be the orthogonal of \(G'\) in E, so that \(\dim(E/F) = p\) . It follows from no. 1, Theorem 2 that there exist p elements \(s_{i} \in S\) such that the \(e_{s_{i}} (1 \leqslant i \leqslant p)\) form a basis of a supplement F of E (applying Theorem 2 (no. 1) to a free subset consisting of a basis of E and the generating system the union of this free subset and the canonical basis of E); we then take the \(f_{i}\) to be the elements of a basis of G' dual to the basis of E/F consisting of the classes of the \(e_{s_{i}} \mod F\) .

COROLLARY 4. Let E be a vector space and M, N two subspaces of E of finite co-dimension; if M', N' are the orthogonals of M and N in E*, the orthogonal of M ∩ N in E* is M' + N'.

As M (resp. N) is the orthogonal of \(M'\) (resp. \(N'\) ) in E (Theorem 7), \(M \cap N\) is the orthogonal of \(M' + N'\) in E and hence \(M' + N'\) is the orthogonal of \(M \cap N\) in \(E^*\) (Theorem 7 (iii)).

COROLLARY 5. Let E be a vector space of finite dimension n.~For every subspace F of E of dimension p, the orthogonal F' of F in E* is of dimension n - p.~For every subspace G' of E* of dimension q, the orthogonal G of G' in E is of dimension n - q and G' is the orthogonal of G in E*.

Theorem 7 gives another characterization of hyperplanes in E:

PROPOSITION 11. For every hyperplane H in a vector space E, there exists a linear form \(x_0^*\) on E such that \(\mathrm{H} = x_0^{*-1}(0)\) . Given such a form \(x_0^*\) , for a linear form \(x^*\) on E to satisfy \(\mathrm{H} = x^{*-1}(0)\) , it is necessary and sufficient that \(x^* = x_0^*\alpha\) , where \(\alpha\) is a scalar \(\neq 0\) . Conversely, for every linear form \(x^* \neq 0\) on E, the subspace \(x^{*-1}(0)\) is a hyperplane of E.

This statement merely expresses Theorem 7 for subspaces of E of codimension 1 and subspaces of E* of dimension 1.

If H is a hyperplane and \(x_{0}^{*}\) a linear form such that \(\mathrm{H}=x^{*-1}(0)\) , the relation

\[
\langle x, x _ {0} ^ {*} \rangle = 0
\]

which characterizes the elements \(x \in \mathbf{H}\) , is called an equation of \(\mathbf{H}\) .

More generally, if \((x_{\iota}^{*})\) is a family of linear forms on E and F denotes the vector subspace the intersection of the hyperplanes \(x_{\iota}^{*-1}(0)\) , the relation ``for all \(\iota\) , \(\langle x, x_{\iota}^{*}\rangle = 0\) '' characterizes the elements of F; the relations

\[
\langle x, x _ {\iota} ^ {*} \rangle = 0 \quad \text { for   all } \iota
\]

form a system of equations of the subspace F. Theorem 7 (ii) expresses the fact that every vector subspace of E can be defined by a system of equations.

Theorem 7 (i) and (ii) proves moreover that a subspace F of finite codimension p can be defined by a system of p equations

\[
\langle x, x _ {i} ^ {*} \rangle = 0, \quad 1 \leqslant i \leqslant p,\tag{19}
\]

where the forms \(x_{i}^{*}\) are linearly independent. Conversely, Corollary 1 to Theorem 7 shows that a subspace \(F\) defined by a system of \(p\) equations (19) is of codimension \(\leqslant p\) and that it is of codimension \(p\) if and only if the \(x_{i}^{*}\) are linearly independent; it amounts to the same to say that F cannot be defined by a system consisting of at most p - 1 of the equations (19).

\subsection{LINEAR EQUATIONS IN VECTOR SPACES}

PROPOSITION 12. Let E, F be two vector spaces over a field K and u:E → F a linear mapping. For the linear equation

\[
u (x) = y _ {0}\tag{20}
\]

to have at least one solution \(x \in \mathbf{E}\) , it is necessary and sufficient that \(y_0\) be orthogonal to the kernel of the transpose mapping \(^t u\) .

The orthogonal of \(u(\mathbf{E})\) in \(\mathbf{F}^*\) is \({}^t u^{-1}(0)\) (§ 2, no. 5, Corollary to Proposition 8) and the orthogonal of \({}^t u^{-1}(0)\) in \(\mathbf{F}\) is therefore \(u(\mathbf{E})\) (no. 5, Theorem 7 (ii)).

We shall obtain a more convenient criterion for systems of scalar linear equations

\[
\langle x, x _ {\iota} ^ {*} \rangle = \eta_ {\iota} \quad (\iota \in I)\tag{21}
\]

where the unknown x takes its values in a vector space E over a field K, the \(x_{\iota}^{*}\) are linear forms on E and the right hand sides \(\eta_{\iota}\) elements of K.

If we consider a basis \((a_{\lambda})_{\lambda\in L}\) of E, the system (21) is equivalent to the system of equations

\[
\sum_ {\lambda \in L} \xi_ {\lambda} \left\langle a _ {\lambda}, x _ {\iota} ^ {*} \right\rangle = \eta_ {\iota} \quad (\iota \in I)\tag{22}
\]

with \(x = \sum_{\lambda \in L} \xi_{\lambda} a_{\lambda}\) , the solutions of (22) being of necessity families \((\xi_{\lambda})\) of elements of \(K\) of finite support.

DEFINITION 4. The dimension of the subspace of \(\mathbf{E}^*\) generated by the family \((x_{\iota}^{*})\) is called the rank of the system (21).

PROPOSITION 13. For the system (21) to be of finite rank \(r\) , it is necessary and sufficient that the linear mapping \(u: x \mapsto (\langle x, x_{\iota}^* \rangle)\) of \(\mathbf{E}\) into \(\mathbf{K}_s^I\) be of rank \(r\) .

If \(F'\) is the subspace of \(E^{*}\) generated by the \(x_{\iota}^{*}\) , the kernel of u is the orthogonal F of \(F'\) in E; if \(F'\) is of dimension r, F is of codimension r and conversely (no. 5, Theorem 7) and \(\operatorname{rg}(u) = \operatorname{codim}_{\mathbf{E}} \mathbf{F}\) (no. 4, formula (11)).

THEOREM 8. Let

\[
\langle x, x _ {\iota} ^ {*} \rangle = \eta_ {\iota} \quad (\iota \in I)\tag{21}
\]

be a system of scalar linear equations on a vector space E over a field K. For this system to have at least one solution, it is necessary that, for every family \((\rho_{\iota})\) of scalars of finite support such that \(\sum_{\iota} x_{\iota}^{*}\rho_{\iota}=0\) , \(\sum_{\iota}\eta_{\iota}\rho_{\iota}=0\) . If the rank of the system (21) is finite, this condition is also sufficient.

The condition is obviously necessary. It says that, if \(F'\) is the subspace of \(E^{*}\) generated by the family \((x_{i}^{*})\) , there exists a linear mapping \(f:F'\to K_{d}\) such that \(f(x_{\iota}^{*})=\eta_{\iota}\) for all \(\iota\in I\) . If \(F'\) is of finite dimension r, \(F'\) is the orthogonal of a subspace F of E of codimension r (no. 5, Theorem 7) and \(F'\) is identified with the dual of E/F (§2, no. 6, Corollary to Proposition 9); f is therefore an element of the bidual \((\mathrm{E}/\mathrm{F})^{**}\) . As E/F is finite-dimensional there exists one and only one element \(y\in E/F\) such that \(f(x^{*})=\langle y,x^{*}\rangle\) for all \(x^{*}\in F'\) (no. 5, Theorem 6). The solutions of (21) are then the \(x\in E\) whose canonical image in E/F is y.

Remark. When the rank of the system (21) is infinite, the condition of Theorem 8 is no longer sufficient. For example, suppose that the \(x_{\iota}^{*}\) are the coordinate forms on the space \(\mathbf{E} = \mathbf{K}_{s}^{(\mathrm{I})}\) , I being infinite (§ 2, no. 6); as the \(x_{\iota}^{*}\) are linearly independent, the condition of Theorem 8 holds for every family ( \(\eta_{\iota}\) ) but the system (21) only has solutions if the family ( \(\eta_{\iota}\) ) has finite support.

A system (21) is always of finite rank if it has only a finite number of equations and its rank is then at most equal to the number of equations (no. 2, Proposition 1). Similarly, if E is of finite dimension n (which for a system (22) corresponds to the case where there is only a finite number n of unknowns), its dual \(E^{*}\) is of dimension n and hence the rank of system (21) is at most equal to n (no. 3, Corollary 4 to Proposition 4). From this we deduce:

COROLLARY 1. A system of scalar linear equations in a vector space, consisting of a finite number of equations whose left hand sides are linearly independent forms, always admits solutions.

COROLLARY 2. For a homogeneous linear system (22) of equations in n unknowns with coefficients in a field K to admit non-trivial solutions consisting of elements of K, it is necessary and sufficient that its rank be \textless n.

This will always be the case if the number of equations is finite and \textless n.

COROLLARY 3. For a linear system (22) with coefficients and right hand sides in a field K, consisting of n equations in n unknowns, to have one and only one solution consisting of elements of K, it is necessary and sufficient that the associated homogeneous system have no non-trivial solution (or, what amounts to the same, that the left hand sides of this system be linearly independent forms).

\subsection{TENSOR PRODUCTS OF VECTOR SPACES}

The results of §§ 3, 4 and 5 relating to free or projective modules apply in particular to vector spaces and give the following properties:

PROPOSITION 14. Given an exact sequence

\[
0 \rightarrow \mathbf {E} ^ {\prime} \rightarrow \mathbf {E} \rightarrow \mathbf {E} ^ {\prime \prime} \rightarrow 0\tag{23}
\]

of right vector spaces over a field K and linear mappings and a left vector space F over K, the corresponding sequence of Z-linear mappings

\[
0 \rightarrow \mathrm{E} ^ {\prime} \otimes_ {\mathbb {K}} \mathrm{F} \rightarrow \mathrm{E} \otimes_ {\mathbb {K}} \mathrm{F} \rightarrow \mathrm{E} ^ {\prime \prime} \otimes_ {\mathbb {K}} \mathrm{F} \rightarrow 0
\]

is exact and splits.

As the sequence (23) splits, this is a particular case of § 3, no. 7, Corollary 5 to Proposition 7 and § 3, no. 6, Proposition 5.

Because of Proposition 14, when \(\mathbf{E}'\) is a vector subspace of \(\mathbf{E}\) , \(j: \mathbf{E}' \to \mathbf{E}\) the canonical injection, \(\mathbf{E}' \otimes_{\mathbb{K}} \mathbf{F}\) is usually identified with a sub- \(\mathbf{Z}\) -module of \(\mathbf{E} \otimes_{\mathbb{K}} \mathbf{F}\) by means of the injection \(j \otimes 1_{\mathbf{F}}\) . With this convention:

COROLLARY. Let K be a field, E a right vector space over K, F a left vector space over K, \((\mathbf{M}_{\alpha})_{\alpha \in \mathbf{A}}\) a family of vector subspaces of E and \((\mathbf{N}_{\beta})_{\beta \in \mathbf{B}}\) a family of vector subspaces of F. Then

\[
\left(\bigcap_ {\alpha \in \mathrm{A}} \mathrm{M} _ {\alpha}\right) \otimes_ {\mathrm{K}} \left(\bigcap_ {\beta \in \mathrm{B}} \mathrm{N} _ {\beta}\right) = \bigcap_ {(\alpha , \beta) \in \mathrm{A} \times \mathrm{B}} \left(\mathrm{M} _ {\alpha} \otimes_ {\mathrm{K}} \mathrm{N} _ {\beta}\right).\tag{24}
\]

It is obviously sufficient to prove the particular case

\[
\left(\bigcap_ {\alpha \in \mathrm{A}} \mathrm{M} _ {\alpha}\right) \otimes_ {\mathrm{K}} \mathrm{F} = \bigcap_ {\alpha \in \mathrm{A}} \left(\mathrm{M} _ {\alpha} \otimes_ {\mathrm{K}} \mathrm{F}\right).\tag{25}
\]

Clearly the left hand side of (25) is contained in the right hand side. To prove the converse, we consider a basis \((f_{\lambda})_{\lambda \in L}\) of F. Every element of \(\mathbf{E} \otimes_{\mathbb{K}} \mathbf{F}\) can then be expressed uniquely in the form \(\sum_{\lambda \in L} x_{\lambda} \otimes f_{\lambda}\) , where \(x_{\lambda} \in \mathbf{E}\) (§3, no. 7, Corollary 1 to Proposition 7); if \(\mathbf{E}'\) is a vector subspace of E, the relation \(\sum_{\lambda \in L} x_{\lambda} \otimes f_{\lambda} \in \mathbf{E}' \otimes_{\mathbb{K}} \mathbf{F}\) is equivalent, by Proposition 14, to \(x_{\lambda} \in \mathbf{E}'\) for all \(\lambda \in L\) . To say that \(\sum_{\lambda \in L} x_{\lambda} \otimes f_{\lambda}\) belongs to each of the \(\mathbf{M}_{\alpha} \otimes_{\mathbb{K}} \mathbf{F}\) thus means that for all \(\lambda \in L\) and all \(\alpha \in A\) , \(x_{\lambda} \in M_{\alpha}\) , that is \(x_{\lambda} \in \bigcap_{\alpha \in A} M_{\alpha}\) for all \(\lambda \in L\) , which proves that the right hand side of (25) is contained in the left hand side. PROPOSITION 15. If \((E_{\lambda})_{\lambda \in L}\) is a family of right vector spaces over a field K and \((F_{\mu})_{\mu \in M}\) a family of left vector spaces over K, the canonical mapping

\[
\left(\prod_ {\lambda \in L} E _ {\lambda}\right) \otimes_ {K} \left(\prod_ {\mu \in M} F _ {\mu}\right)\rightarrow \prod_ {(\lambda , \mu) \in L \times M} (E _ {\lambda} \otimes_ {K} F _ {\mu})\tag{26}
\]

(§ 3, no. 7, formula (22)) is injective.

We write \(\mathbf{F} = \prod_{\mu \in \mathbf{M}} \mathbf{F}_{\mu}\) ; the mapping (26) is the composition of the canonical mappings

\[
\left(\prod_ {\lambda \in L} \mathbf {E} _ {\lambda}\right) \otimes_ {\mathbb {K}} F \rightarrow \prod_ {\lambda \in L} \left(\mathbf {E} _ {\lambda} \otimes_ {\mathbb {K}} F\right)
\]

and

\[
\prod_ {\lambda \in L} \left(\mathrm{E} _ {\lambda} \otimes_ {K} F\right)\rightarrow \prod_ {\lambda \in L} \left(\prod_ {\mu \in M} \left(\mathrm{E} _ {\lambda} \otimes_ {K} F _ {\mu}\right)\right);
\]

as F and the \(\mathbf{E}_{\lambda}\) are vector spaces over K, this reduces to §3, no. 7, Corollary 3 to Proposition 7.

When the conditions of Proposition 15 are fulfilled, the tensor product \(\left(\prod_{\lambda \in L} E_{\lambda}\right) \otimes_{K} \left(\prod_{\mu \in M} F_{\mu}\right)\) is often identified with its canonical image in

\[
\prod_ {\lambda , \mu} (\mathbf {E} _ {\lambda} \otimes_ {K} \mathbf {F} _ {\mu}).
\]

With this convention:

COROLLARY. Let F be a left vector space over K; for every set X, the left vector space \(K_{d}^{X} \otimes_{K} F\) is identified with the subspace of the space \(F^{X}\) of all mappings of X into F, consisting of the mappings u such that \(u(X)\) is of finite rank in F.

If \((f_{\lambda})\) is a basis of \(F\) , the element \(\sum_{\lambda \in L} v_{\lambda} \otimes f_{\lambda}\) of \(K_d^X \otimes_K F\) is identified under (26) with the mapping \(x \mapsto \sum_{\lambda} v_{\lambda}(x)f_{\lambda}\) . As \(v_{\lambda} = 0\) except for the indices \(\lambda\) belonging to a finite subset \(H\) of \(L\) , the image of \(X\) under the above mapping is contained in the finite-dimensional subspace of \(F\) generated by the \(f_{\lambda}\) of indices \(\lambda \in H\) . Conversely, let \(u: X \to F\) be a mapping such that \(u(X)\) is contained in a finite-dimensional subspace \(G\) of \(F\) and let \((b_i)_{1 \leqslant i \leqslant n}\) be a basis of \(G\) . For all \(x \in X\) , we may write \(u(x) = \sum_{i=1}^{n} v_i(x)b_i\) , where the \(v_i(x)\) are well determined elements of \(K\) ; thus \(n\) mappings \(v_i: X \to K\) are defined and clearly \(u\) is then identified with the element \(\sum_{i=1}^{n} v_i \otimes b_i\) .

Similarly, for a right vector space E over K and a set Y, \(E \otimes_{K} K_{s}^{Y}\) is identified with a subspace of the space \(E^{Y}\) , consisting of the mappings \(v: Y \to E\) such that \(v(Y)\) is of finite rank. More particularly, for every field K, \(K_{d}^{X} \otimes_{K} K_{s}^{Y}\) is identified with a subspace of the space \(K^{X \times Y}\) of mappings of \(X \times Y\) into K (K being considered as a (K, K)-bimodule); an element \(\sum_{i} u_{i} \otimes v_{i}\) , where \(u_{i}\) is a mapping of X into K and \(v_{i}\) a mapping of Y into K, is identified with the mapping \((x, y) \mapsto \sum_{i} u_{i}(x) v_{i}(y)\) of \(X \times Y\) into K.

PROPOSITION 16. (i) Let K, L be two fields, E a left vector space over K, F a left vector space over L and G a (K, L)-bimodule. Then the canonical Z-homomorphism

\[
\nu : \operatorname{Hom} _ {\mathbf {K}} (\mathrm{E}, \mathrm{G}) \otimes_ {\mathrm{L}} \mathrm{F} \rightarrow \operatorname{Hom} _ {\mathbf {K}} (\mathrm{E}, \mathrm{G} \otimes_ {\mathrm{L}} \mathrm{F})\tag{27}
\]

(§ 4, no. 2, formula (7)) is injective; it is bijective when one of the vector spaces E, F is finite-dimensional.

\begin{enumerate}
\def\labelenumi{(\roman{enumi})}
\setcounter{enumi}{1}
\tightlist
\item
  Let \(E_{1}\) , \(E_{2}\) , \(F_{1}\) , \(F_{2}\) be four vector spaces over a commutative field K; then the canonical K-homomorphism
\end{enumerate}

\[
\lambda : \operatorname{Hom} \left(\mathrm{E} _ {1}, \mathrm{F} _ {1}\right) \otimes \operatorname{Hom} \left(\mathrm{E} _ {2}, \mathrm{F} _ {2}\right)\rightarrow \operatorname{Hom} \left(\mathrm{E} _ {1} \otimes \mathrm{E} _ {2}, \mathrm{F} _ {1} \otimes \mathrm{F} _ {2}\right)\tag{28}
\]

(§ 4, no. 4, formula (21)) is injective; it is bijective if one of the ordered pairs \((\mathbf{E}_1, \mathbf{E}_2), (\mathbf{E}_1, \mathbf{F}_1), (\mathbf{E}_2, \mathbf{F}_2)\) consists of finite-dimensional spaces.

Assertion (i) is a particular case of § 4, no. 2, Proposition 2. Similarly the second assertion of (ii) is a particular case of § 4, no. 4, Proposition 4. Finally, to see that the homomorphism (28) is always injective, observe that \(\operatorname{Hom}(\mathbf{E}_i, \mathbf{F}_i)\) is a vector subspace of \(\mathbf{F}_i^{\mathbf{E}_i}\) ( \(i = 1, 2\) ) and that

\[
\operatorname{Hom} \left(\mathrm{E} _ {1} \otimes \mathrm{E} _ {2}, \mathrm{F} _ {1} \otimes \mathrm{F} _ {2}\right)
\]

is canonically identified with a vector subspace of the space \((\mathbf{F}_{1} \otimes \mathbf{F}_{2})^{\mathbf{E}_{1} \times \mathbf{E}_{2}}\) (II, §3, no. 1, Proposition 1); when these identifications are made and also the left hand side of (28) is identified with a subspace of \(F_{1}^{E_{1}} \otimes F_{2}^{E_{2}}\) (Proposition 14), the canonical mapping (28) becomes the restriction to this subspace of the canonical mapping (26) and it has been seen (Proposition 15) that the latter is injective.

COROLLARY 1. Let E and F be two vector spaces over a field K; the canonical mapping

\[
\mathbf {E} ^ {*} \otimes_ {\mathbb {K}} \mathbf {F} \rightarrow \operatorname{Hom} _ {\mathbb {K}} (\mathbf {E}, \mathbf {F})
\]

(§ 4, no. 2, formula (11)) is injective; it is bijective when E or F is finite-dimensional.

This is a special case of Proposition 16 (i).

COROLLARY 2. Let E be a right vector space and F a left vector space over the same field K; the canonical mapping

\[
\mathbf {E} \otimes_ {\mathbb {K}} \mathbf {F} \rightarrow \operatorname{Hom} _ {\mathbb {K}} (\mathbf {E} ^ {*}, \mathbf {F})\tag{29}
\]

(§ 4, no. 2, formula (15)) is injective; it is bijective when E is finite-dimensional.

This is a special case of § 4, no. 2, Remark 2.

Remarks. (1) Let K be a commutative field, E, F two vector spaces over K,

\((a_{\lambda})\) a basis of \(\mathbf{E}\) and \((b_{\mu})\) a basis of \(\mathrm{F}\) ; then \((a_{\lambda} \otimes b_{\mu})\) is a basis of the vector \(\mathbf{K}\) -space \(\mathbf{E} \otimes_{\mathbb{K}} \mathbf{F}\) (II, § 3, no. 7, Corollary 2 to Proposition 2) and therefore

\[
\dim_ {\mathbf {K}} (\mathbf {E} \otimes_ {\mathbf {K}} \mathbf {F}) = \dim_ {\mathbf {K}} \mathbf {E}. \dim_ {\mathbf {K}} \mathbf {F}.\tag{30}
\]

\begin{enumerate}
\def\labelenumi{(\arabic{enumi})}
\setcounter{enumi}{1}
\tightlist
\item
  Let K be a commutative field, \(E_{1}\) , \(E_{2}\) , \(F_{1}\) , \(F_{2}\) four vector spaces over K and \(u: E_{1} \rightarrow F_{1}\) , \(v: E_{2} \rightarrow F_{2}\) two linear mappings; then
\end{enumerate}

\[
\operatorname{rg} (u \otimes v) = \operatorname{rg} (u). \operatorname{rg} (v).\tag{31}
\]

It is immediate that \((u\otimes v)(\mathbf{E}_1\otimes \mathbf{E}_2)\) is the canonical image of \(u(\mathbf{E}_1)\otimes v(\mathbf{E}_2)\) in \(\mathbf{F}_1\otimes \mathbf{F}_2\) and hence (Proposition 14) is isomorphic to \(u(\mathbf{E}_1)\otimes v(\mathbf{E}_2)\) ; the conclusion then follows from (30).

\begin{enumerate}
\def\labelenumi{(\arabic{enumi})}
\setcounter{enumi}{2}
\tightlist
\item
  Under the same hypotheses as in Remark 1,
\end{enumerate}

\[
\dim_ {\mathbb {K}} (\operatorname{Hom} _ {\mathbb {K}} (\mathrm{E}, \mathrm{F})) \geqslant \dim_ {\mathbb {K}} \mathrm{E}. \dim_ {\mathbb {K}} \mathrm{F}.\tag{32}
\]

If E is isomorphic to \(K^{(I)}\) , \(\operatorname{Hom}(E, F)\) is isomorphic to \((\operatorname{Hom}(K, F))^{I}\) (§ 1, no. 6, Corollary 1 to Proposition 6) and hence to \(F^{I}\) (§ 1, no. 14); as \(F^{(I)}\) is a subspace of \(F^{I}\) and \(\dim(F^{(I)}) = \operatorname{Card}(I).dim F = \dim E. dim F\) (no. 2, Proposition 2), the inequality (32) follows from no. 3, Corollary 4 to Proposition 4. The same argument shows that the two sides of (32) are equal when dim E is finite (cf.~§ 10, nos. 3 and 4).

\subsection{RANK OF AN ELEMENT OF A TENSOR PRODUCT}

Let E be a right vector space and F a left vector space over the same field K; to every element \(u \in E \otimes_{K} F\) there corresponds canonically under (29) a homomorphism \(u_{1} \in \operatorname{Hom}_{\mathbb{K}}(\mathbf{E}^{*}, \mathbf{F})\) ; if \(u = \sum_{i} x_{i} \otimes y_{i}\) with \(x_{i} \in E, y_{i} \in F\) , the element \(u_{1}\) is the linear mapping

\[
x ^ {*} \mapsto \sum_ {i} \langle x ^ {*}, x _ {i} \rangle y _ {i}.\tag{33}
\]

On the other hand, \(\mathbf{E} \otimes_{\mathbb{K}} \mathbf{F}\) is canonically identified with \(\mathbf{F} \otimes_{\mathbb{K}^0} \mathbf{E}\) , where \(\mathbf{E}\) is considered as a left vector space and \(\mathbf{F}\) as a right vector space over the opposite field \(\mathbf{K}^0\) ; thus there corresponds canonically to \(u\) a homomorphism \(u_2 \in \operatorname{Hom}_{\mathbb{K}}(\mathbf{F}^*, \mathbf{E})\) given by

\[
y ^ {*} \mapsto \sum_ {i} x _ {i} \langle y _ {i}, y ^ {*} \rangle ;\tag{34}
\]

\(u_{1}\) (resp. \(u_{2}\) ), considered as a mapping of \(E^{*}\) into \(F^{**}\) (resp. of \(F^{*}\) into \(E^{**}\) ) is just the transpose of \(u_{2}\) (resp. \(u_{1}\) ). The ranks of \(u_{1}\) and \(u_{2}\) are thus equal to the same finite number r, the common dimension of the subspaces \(u_{1}(E^{*})\) of F and \(u_{2}(F^{*})\) of E, each of which is canonically isomorphic to the dual of the other (no. 5); we shall say that r (denoted \(\operatorname{rg}(u)\) ) is the rank of the element u of \(E \otimes_{K} F\) and that \(u_{1}(E^{*})\) and \(u_{2}(F^{*})\) are the subspaces (of F and E respectively) associated with u.

PROPOSITION 17. Let \(u\) be an element of \(\mathbf{E} \otimes_{\mathbb{K}} \mathbf{F}\) and \(\mathbf{M} \subset \mathbf{E}\) and \(\mathbf{N} \subset \mathbf{F}\) its associated subspaces. For every expression \(u = \sum_{i=1}^{s} x_i \otimes y_i\) of \(u\) , where \(x_i \in \mathbf{E}\) and \(y_i \in \mathbf{F}\) for \(1 \leqslant i \leqslant s\) , the subspace \(\mathbf{M}\) (resp. \(\mathbf{N}\) ) is contained in the subspace of \(\mathbf{E}\) (resp. \(\mathbf{F}\) ) generated by the \(x_i\) (resp. the \(y_i\) ). Moreover, the following properties are equivalent:

\begin{enumerate}
\def\labelenumi{(\alph{enumi})}
\item
  The integer \(s\) is equal to the rank of \(u\) .
\item
  The family \((x_{i})_{1\leqslant i\leqslant s}\) is a basis of M.
\item
  The family \((y_{i})_{1\leqslant i\leqslant s}\) is a basis of N.
\item
  The families \((x_{i})_{1\leqslant i\leqslant s}\) and \((y_{i})_{1\leqslant i\leqslant s}\) are both free.
\end{enumerate}

By (33) (resp. (34)) each element of \(N = u_{1}(E^{*})\) (resp. \(M = u_{2}(F^{*})\) ) is a linear combination of the \(y_{i}\) (resp. the \(x_{i}\) ); whence the first assertion. If s = r, the subspace generated by the \(x_{i}\) (resp. \(y_{i}\) ) with dimension \(\leqslant \dim M\) (resp. \(\leqslant \dim N\) ) and containing M (resp. N) is identical with it and hence (a) implies (b) and (c) and a fortiori (d). Conversely, each of conditions (b) and (c) implies (a) by definition of \(\operatorname{rg}(u)\) . Finally if (d) holds, there exists a family \((x_{i}^{*})_{1\leqslant i\leqslant s}\) of elements of \(E^{*}\) such that \(\langle x_{i}, x_{j}^{*}\rangle = \delta_{ij}\) (no. 5, Corollary 1 to Theorem 7) and hence it follows from (33) that \((y_{i})\) is a basis of N, which completes the proof.

COROLLARY 1. The rank of \(u\) is the smallest integer \(s\) such that there exists an expression \(u = \sum_{i=1}^{s} x_i \otimes y_i\) , where \(x_i \in \mathbf{E}\) and \(y_i \in \mathbf{F}\) for \(1 \leqslant i \leqslant s\) .

This follows immediately from Proposition 17 and no. 2, Proposition 1.

COROLLARY 2. Let K be a commutative field, E, F two vector spaces over K and L a commutative extension field of K. Let u be an element of E ⊗K F, M and N its associated subspaces and u' the canonical image of u in (E ⊗K F)(L) (canonically identified with E(L) ⊗L F(L), cf.~§ 5, no. 1, Proposition 3); then rg(u') = rg(u) and the subspaces associated with u' are canonically identified with M(L) and N(L).

If \(u = \sum_{i=1}^{r} x_i \otimes y_i\) , where the families \((x_i)\) and \((y_i)\) are free, then \(u' = \sum_{i=1}^{r} (1 \otimes x_i) \otimes (1 \otimes y_i)\) and the families \((1 \otimes x_i)\) and \((1 \otimes y_i)\) are free in \(E_{(L)}\) and \(F_{(L)}\) respectively (§ 5, no. 1, Proposition 4).

\subsection{EXTENSION OF SCALARS FOR A VECTOR SPACE}

Recall (I, § 9, no. 1, Theorem 2) that a homomorphism of a field K into a non-zero ring A is necessarily injective.

PROPOSITION 18. Let \(\rho\) be a homomorphism of a field \(K\) into a ring \(A\) . For every exact sequence of vector \(K\) -spaces and \(K\) -linear mappings

\[
\mathrm{E} ^ {\prime} \xrightarrow {u} \mathrm{E} \xrightarrow {v} \mathrm{E} ^ {\prime \prime}
\]

the sequence

\[
\mathrm{E} _ {(\mathrm{A})} ^ {\prime} \xrightarrow {u _ {(\mathrm{A})}} \mathrm{E} _ {(\mathrm{A})} \xrightarrow {v _ {(\mathrm{A})}} \mathrm{E} _ {(\mathrm{A})} ^ {\prime \prime}
\]

is exact.

This is a special case of no. 7, Proposition 14, taking account of § 1, no. 4, Remark 4.

COROLLARY. For every K-linear mapping \(f: \mathrm{E}' \to \mathrm{E}\) , \(\operatorname{Im}(f_{(\mathbf{A})}) = (\operatorname{Im}(f))_{(\mathbf{A})}\) , \(\operatorname{Ker}(f_{(\mathbf{A})}) = (\operatorname{Ker}(f))_{(\mathbf{A})}\) , \(\operatorname{Coker}(f_{(\mathbf{A})}) = (\operatorname{Coker}(f))_{(\mathbf{A})}\) , to within canonical isomorphisms.

PROPOSITION 19. Let \(\rho\) be an injective homomorphism of a field K into a ring A. For every left vector space E over K, the canonical mapping \(\phi: \mathbf{E} \to \rho^{*}(\mathbf{E}) = \mathbf{A} \otimes_{\mathbb{K}} \mathbf{E}\) is injective. Moreover, for every vector subspace \(\mathbf{E}'\) of E, \(\rho^{*}(\mathbf{E}') = \mathbf{A} \otimes_{\mathbb{K}} \mathbf{E}'\) is canonically identified with a direct factor sub-A-module of \(\mathbf{A} \otimes_{\mathbb{K}} \mathbf{E}\) and, with this identification,

\[
(\mathbf {A} \otimes_ {\mathbf {K}} \mathbf {E} ^ {\prime}) \cap \phi (\mathbf {E}) = \phi (\mathbf {E} ^ {\prime}).\tag{35}
\]

The first assertion is a special case of § 5, no. 1, Proposition 4; the second is a special case of no. 7, Proposition 14; finally, to show (35), it suffices to take in A (considered as a right K-module) a basis \((a_{\lambda})_{\lambda \in L}\) such that \(a_{\lambda_0} = 1\) for some index \(\lambda_0\) (no. 1, Theorem 2); the elements of A \(\otimes_{\mathbb{K}}\) E can be written uniquely as \(\sum_{\lambda} a_{\lambda} \otimes x_{\lambda}\) with \(x_{\lambda} \in \mathbf{E}\) and, for such an element to belong to A \(\otimes_{\mathbb{K}}\) E', it is necessary and sufficient that \(x_{\lambda} \in \mathbf{E}'\) for all \(\lambda\) . On the other hand, the elements of \(\phi(\mathbf{E})\) are those for which \(x_{\lambda} = 0\) for \(\lambda \neq \lambda_0\) ; for an element \(\sum_{\lambda} a_{\lambda} \otimes x_{\lambda}\) to belong to (A \(\otimes_{\mathbb{K}}\) E') \(\cap\) \(\phi(\mathbf{E})\) , it is necessary and sufficient that \(x_{\lambda} = 0\) for \(\lambda \neq \lambda_0\) and \(x_{\lambda_0} \in \mathbf{E}'\) , whence the conclusion.

COROLLARY. Let \(\rho\) be an injective homomorphism of a field K into a ring A. For a K-linear mapping \(f: \mathbf{E} \to \mathbf{F}\) (where E and F are two vector spaces over K) to be injective (resp. surjective, zero), it is necessary and sufficient that \(f_{(\mathbf{A})}: \mathbf{E}_{(\mathbf{A})} \to \mathbf{F}_{(\mathbf{A})}\) be injective (resp. surjective, zero).

This follows immediately from Proposition 19 and the Corollary to Proposition 18.

PROPOSITION 20. Let \(\rho\) be a homomorphism of a field \(\mathbf{K}\) into a ring \(\mathbf{A}\) . For every left vector space \(\mathbf{E}\) over \(\mathbf{K}\) , the canonical right A-module homomorphism

\[
\upsilon : \left(\mathrm{E} ^ {*}\right) _ {(\mathrm{A})} \rightarrow \left(\mathrm{E} _ {(\mathrm{A})}\right) ^ {*}
\]

(§ 5, no. 4) is injective; it is bijective when E is finite-dimensional.

The second assertion follows from § 5, no. 4, Proposition 8. To prove the first, we note that every element of \((\mathbf{E}^{*})_{(\mathbf{A})}\) may be written as \(\sum_{i} x_{i}^{*} \otimes \alpha_{i}\) , where \(\alpha_{i} \in \mathbf{A}\) and \((x_{i}^{*})_{1 \leqslant i \leqslant n}\) is a free family in \(\mathbf{E}^{*}\) ; there corresponds to it in \((\mathbf{E}_{(\mathbf{A})})^{*}\) the linear form \(y^{*}\) such that \(y^{*}(1 \otimes x) = \sum_{i} \rho(\langle x, x_{i}^{*} \rangle)\alpha_{i}\) for all \(x \in \mathbf{E}\) . But, there exists in E a family \((x_{i})_{1 \leqslant i \leqslant n}\) such that \(\langle x_{i}, x_{j}^{*} \rangle = \delta_{ij}\) (no. 5, Corollary 2 to Theorem 7), whence \(y^{*}(1 \otimes x_{i}) = \alpha_{i}\) ; the relation \(y^{*} = 0\) therefore implies \(\alpha_{i} = 0\) for all \(i\) , which proves our assertion.

PROPOSITION 21. Let K be a field and L an extension field of K.

\begin{enumerate}
\def\labelenumi{(\roman{enumi})}
\item
  For every vector space \(\mathbf{E}\) over \(\mathbf{K}\) , \(\dim_{\mathbf{L}}(\mathrm{E}_{(\mathbf{L})}) = \dim_{\mathbf{K}}\mathbf{E}\) .
\item
  For every K-linear mapping \(u: \mathbf{E} \to \mathbf{F}\) , where \(\mathbf{E}\) and \(\mathbf{F}\) are vector spaces over \(\mathbf{K}\) , \(\operatorname{rg}(u_{(L)}) = \operatorname{rg}(u)\) .
\end{enumerate}

If \((e_{\iota})_{\iota\in I}\) is a basis of \(E\) over \(K\) , \((1\otimes e_{\iota})_{\iota\in I}\) is a basis of \(\mathbf{E}_{(L)}\) over \(L\) (§ 5, no. 1, Proposition 4), whence the first assertion; the second follows from the first and the fact that \(u_{(L)}(\mathbf{E}_{(L)})\) is canonically identified with \((u(\mathbf{E}))_{(L)}\) by the Corollary to Proposition 18.

PROPOSITION 22. Let K be a commutative field, \(\rho: \mathrm{K} \to \mathrm{A}\) an injective central homomorphism and E, F two vector spaces over K. Then the canonical homomorphism

\[
\omega : \mathrm{A} \otimes_ {\mathrm{K}} \operatorname{Hom} (\mathrm{E}, \mathrm{F}) \rightarrow \operatorname{Hom} _ {\mathrm{A}} \left(\mathrm{E} _ {(\mathrm{A})}, \mathrm{F} _ {(\mathrm{A})}\right)\tag{36}
\]

(§ 5, no. 3, formula (17)) is injective; it is bijective if A or E is a finite-dimensional vector space over K.

This is a particular case of § 5, no. 3, Proposition 7.

\subsection{MODULES OVER INTEGRAL DOMAINS}

PROPOSITION 23. In a module E over an integral domain A, the set T of non-free elements is a submodule of E.

If x and y are not free, there exist two non-zero elements \(\alpha\) , \(\beta\) in A such that \(\alpha x = 0\) and \(\beta y = 0\) . Then \(\alpha\beta \neq 0\) since A is an integral domain and \(\alpha\beta(\lambda x + \mu y) = 0\) for all \(\lambda\) and \(\mu\) in A since A is commutative, hence \(\lambda x + \mu y\) is not free.

Remark. Let E be a module over any commutative ring A. If x is a non-free element of E, every element of the submodule Ax is non-free. On the other hand, if A contains divisors of 0, the sum of two non-free elements of E may be free; for example, in Z/6Z considered as a module over itself, 3 and 4 are not free, but \(3 + 4 = 1\) is free.

Proposition 23 leads to the following definition:

DEFINITION 5. In a module E over an integral domain A, the torsion submodule of E is the submodule of E consisting of the non-free elements (also called the torsion elements of E).

When E is equal to its torsion submodule (that is when every element of E is annihilated by an element \(\neq0\) of A) E is called a torsion module. When the torsion submodule of E is reduced to 0 (that is every non-zero element of E is free) E is called (by an abuse of language) a torsion-free module.

Every submodule of a free A-module (and in particular every projective A-module) is torsion-free. The Z-module Q is torsion-free.

PROPOSITION 24. Let A be an integral domain. For every A-module E, let T(E) denote the torsion submodule of E. Let \(f: \mathbf{E} \to \mathbf{E}'\) be an A-linear mapping, E and \(\mathbf{E}'\) being A-modules.

\begin{enumerate}
\def\labelenumi{(\roman{enumi})}
\item
  \(f(\mathbf{T}(\mathbf{E})) \subset \mathbf{T}(\mathbf{E}')\) .
\item
  If \(f\) is injective, \(f(\mathbf{T}(\mathbf{E})) = \mathbf{T}(\mathbf{E}') \cap f(\mathbf{E})\) .
\item
  If \(f\) is surjective and \(\operatorname{Ker}(f) \subset \mathrm{T}(\mathbf{E})\) , then \(f(\mathrm{T}(\mathbf{E})) = \mathrm{T}(\mathbf{E}')\) .
\end{enumerate}

Assertions (i) and (ii) are obvious. On the other hand, if \(f\) is surjective and \(x' \in \mathrm{T}(\mathrm{E}')\) , then \(x' = f(x)\) , where \(x \in \mathbf{E}\) , and by hypothesis there exists \(\alpha \neq 0\) in \(\mathbf{A}\) such that \(f(\alpha x) = \alpha x' = 0\) ; whence \(\alpha x \in \operatorname{Ker}(f)\) and by the hypothesis there exists \(\beta \neq 0\) in \(\mathbf{A}\) such that \(\beta(\alpha x) = 0\) ; as \(\beta \alpha \neq 0\) , \(x \in \mathrm{T}(\mathrm{E})\) .

COROLLARY 1. For every A-module E, E/T(E) is torsion-free.

If \(f: \mathbf{E} \to \mathbf{E}'\) is an A-linear mapping, let \(f_{\mathrm{T}}\) denote the mapping \(\mathrm{T}(\mathbf{E}) \to \mathrm{T}(\mathbf{E}')\) with the same graph as the restriction of \(f\) to \(\mathrm{T}(\mathbf{E})\) . With this notation:

COROLLARY 2. For every exact sequence of A-modules and A-linear mappings

\[
0 \longrightarrow \mathrm{E} ^ {\prime} \stackrel {f} {\longrightarrow} \mathrm{E} \stackrel {g} {\longrightarrow} \mathrm{E} ^ {\prime \prime}
\]

the sequence

\[
0 \longrightarrow \mathrm{T} (\mathrm{E} ^ {\prime}) \xrightarrow {f _ {\mathrm{T}}} \mathrm{T} (\mathrm{E}) \xrightarrow {g _ {\mathrm{T}}} \mathrm{T} (\mathrm{E} ^ {\prime \prime})
\]

is exact.

\[
\operatorname{Ker} \left(g _ {\mathrm{T}}\right) = \operatorname{Ker} (g) \cap \mathrm{T} (\mathrm{E}) = f \left(\mathrm{E} ^ {\prime}\right) \cap \mathrm{T} (\mathrm{E}) = f \left(\mathrm{T} \left(\mathrm{E} ^ {\prime}\right)\right) = \operatorname{Im} \left(f _ {\mathrm{T}}\right).
\]

PROPOSITION 25. Let A be an integral domain and \((\mathbf{E}_{\iota})\) a family of A-modules; then

\begin{enumerate}
\def\labelenumi{(\arabic{enumi})}
\setcounter{enumi}{36}
\tightlist
\item
\end{enumerate}

\[
\mathrm{T} \left(\bigoplus_ {\iota} \mathrm{E} _ {\iota}\right) = \bigoplus_ {\iota} \mathrm{T} (\mathrm{E} _ {\iota}).
\]

Let \((x_{\iota})\) be an element of \(\bigoplus_{\iota} E_{\iota}\) such that \(x_{\iota} \in T(E_{\iota})\) for all \(\iota\) ; then each of the \(x_{\iota}\) is annihilated by an element \(\alpha_{\iota} \neq 0\) of A and it may be assumed that \(\alpha_{\iota} = 1\) when \(x_{\iota} = 0\) ; as the family \((x_{\iota})\) has finite support, the element \(\alpha = \prod_{\iota} \alpha_{\iota}\) of A is defined and \(\neq 0\) ; it obviously annihilates \(\bigoplus_{\iota} x_{\iota}\) and hence \(\bigoplus_{\iota} T(E_{\iota}) \subset T\left(\bigoplus_{\iota} E_{\iota}\right)\) ; the converse is immediate.

If E and F are two A-modules, clearly \(T(E \otimes_{A} F)\) contains the canonical images of \(T(E) \otimes_{A} F\) and \(E \otimes_{A} T(F)\) ; but examples can be given of torsion-free A-modules E, F such that \(T(E \otimes_{A} F) \neq 0\) (Exercise 31).

Note that an infinite product of torsion modules is not necessarily a torsion module; for example, in the Z-module \(\prod_{n=1}^{\infty}\left(\mathbf{Z}/p^{n}\mathbf{Z}\right)\) (p an integer \textgreater1), the element all of whose coordinates are 1 is free.

PROPOSITION 26. Let A be an integral domain, K its field of fractions, E an A-module and \(\mathbf{E}_{(\mathbf{K})} = \mathbf{K} \otimes_{\mathbf{A}} \mathbf{E}\) the vector space over K obtained by extending the ring of operators; let \(\phi\) denote the canonical A-linear mapping \(x \mapsto 1 \otimes x\) of E into \(\mathbf{E}_{(\mathbf{K})}\) .

\begin{enumerate}
\def\labelenumi{(\roman{enumi})}
\item
  Every element of \(\mathbf{E}_{(\mathbf{K})}\) is of the form \(\lambda^{-1}\phi(x)\) for \(\lambda \in \mathbf{A}\) , \(\lambda \neq 0\) and \(x \in \mathbf{E}\) .
\item
  The kernel of \(\phi\) is the torsion submodule T(E) of E.
\item
  Every element of \(\mathrm{E}_{(\mathbb{K})}\) is of the form \(z = \sum_{i=1}^{n} \xi_i \phi(x_i)\) with \(\xi_i \in \mathbb{K}\) and \(x_{i}\in \mathbf{E};\) for all \(i\) , there exists \(\alpha_{i}\in \mathbf{A}\) such that \(\alpha_{i}\neq 0\) and \(\alpha_{i}\xi_{i}\in \mathbf{A}\) ; if \(\alpha = \prod_{i = 1}^{n}\alpha_{i},\) then \(\alpha \neq 0\) and \(\alpha \xi_{i} = \beta_{i}\in \mathbf{A}\) for all \(i\) , whence, in \(\mathrm{E}_{(\mathbf{K})}\) ,
\end{enumerate}

\[
z = \alpha^ {- 1} (\alpha z) = \alpha^ {- 1} \sum_ {i = 1} ^ {n} \beta_ {i} \phi (x _ {i}) = \alpha^ {- 1} \phi \left(\sum_ {i = 1} ^ {n} \beta_ {i} x _ {i}\right)
\]

since \(\phi\) is A-linear.

\begin{enumerate}
\def\labelenumi{(\roman{enumi})}
\setcounter{enumi}{1}
\tightlist
\item
  If \(x \neq 0\) is not free in \(\mathbf{E}\) , there exists \(\alpha \neq 0\) in \(\mathbf{A}\) such that \(\alpha x = 0\) , whence \(\alpha \phi(x) = \phi(\alpha x) = 0\) in \(\mathrm{E}_{(\mathbb{K})}\) , which implies \(\phi(x) = 0\) . Conversely, suppose that, for some \(x \in \mathbf{E}\) , \(1 \otimes x = 0\) in \(\mathrm{E}_{(\mathbb{K})}\) ; we show that \(x\) is a torsion element of \(\mathbf{E}\) . We consider the set \(\mathfrak{M}\) of monogenous sub-A-modules of \(\mathbf{K}\) ; this is a right directed set under the relation of inclusion, for any two elements \(\alpha, \beta\) of \(\mathbf{K}\) can be written as \(\alpha = \zeta^{-1}\xi, \beta = \zeta^{-1}\eta\) , where \(\xi, \eta, \zeta\) belong to \(\mathbf{A}\) and \(\zeta \neq 0\) , hence \(\mathbf{A}. \alpha \subset \mathbf{A}. \zeta^{-1}\) and \(\mathbf{A}. \beta \subset \mathbf{A}. \zeta^{-1}\) . Moreover \(\mathbf{K}\) is the union of the modules \(\mathbf{M} \in \mathfrak{M}\) and can therefore be considered as the direct limit of the direct system defined by the modules \(\mathbf{M} \in \mathfrak{M}\) and the canonical injections (§ 6, no. 2, Remark). Hence also, to within a canonical isomorphism, \(\mathrm{E}_{(\mathbb{K})} = \varinjlim (\mathbf{M} \otimes_{\mathbf{A}} \mathbf{E})\) (§ 6, no. 3, Proposition 7) and the relation \(1 \otimes x = 0\) in \(\mathrm{E}_{(\mathbb{K})}\) implies that there exists an
\end{enumerate}

\(M \in \mathfrak{M}\) such that \(1 \in M\) and \(1 \otimes x = 0\) in the tensor product \(M \otimes_{A} E\) (Set Theory, III, § 7, no. 5, Lemma 1). It may further be supposed (replacing if need be \(M\) by a monogenous submodule \(M' \supset M\) of \(K\) ) that \(M = A. \gamma^{-1}\) , where \(\gamma \in A\) and \(\gamma \neq 0\) . Now the mapping \(\xi \mapsto \gamma\xi\) is an isomorphism of \(M\) onto the A-module \(A\) ; on the other hand, the canonical isomorphism \(A \otimes_{A} E \to E\) (§ 3, no. 4, Proposition 4) maps \(\xi \otimes x\) to the element \(\xi x\) of \(E\) ; thus there exists an isomorphism \(M \otimes_{A} E \to E\) which maps the tensor product \(\xi \otimes x\) to the element \((\gamma\xi)x\) of \(E\) . The hypothesis \(1 \otimes x = 0\) in \(M \otimes_{A} E\) thus implies \(\gamma x = 0\) .

Remark. Let \(\alpha^{-1}\phi(x)\) , \(\beta^{-1}\phi(y)\) be two elements of \(\mathbf{E}_{(\mathbb{K})}\) , with \(\alpha\in\mathbf{A}\) , \(\beta\in\mathbf{A}\) , \(x\in\mathbf{E}\) , \(y\in\mathbf{E}\) , \(\alpha\beta\neq0\) . In order that \(\alpha^{-1}\phi(x)=\beta^{-1}\phi(y)\) , it is necessary and sufficient that \(\beta x-\alpha y\) be a torsion element of \(\mathbf{E}\) , for this relation is equivalent to \(\beta\phi(x)=\alpha\phi(y)\) , which may also be written \(\phi(\beta x-\alpha y)=0\) .

COROLLARY 1. If E is a torsion-free A-module, the canonical mapping \(\phi: \mathbf{E} \to \mathbf{E}_{(\mathbb{K})}\) is injective.

Recall (§ 5, no. 1) that for every A-linear mapping of E into a vector space F over K, there exists one and only one K-linear mapping \(\tilde{f}\colon\mathbf{E}_{(\mathbb{K})}\to\mathbf{F}\) such that \(f=\tilde{f}\circ\phi\) ; we shall say that \(\tilde{f}\) is associated with \(f\) .

COROLLARY 2. Let \(f\) be an A-linear mapping of \(\mathbf{E}\) into a vector space \(\mathbf{F}\) over \(\mathbf{K}\) ; if \(\operatorname{Ker}(f) \subset \mathrm{T}(\mathbf{E})\) , the K-linear mapping \(\bar{f}\) associated with \(f\) is injective.

We write an element of \(\operatorname{Ker}(\bar{f})\) in the form \(\lambda^{-1}\phi(x)\) , where \(\lambda \in \mathbf{A}\) , \(\lambda \neq 0\) , \(x \in \mathbf{E}\) ; the relation \(\bar{f}(\lambda^{-1}\phi(x)) = 0\) is equivalent to \(\lambda^{-1}\bar{f}(\phi(x)) = 0\) in \(\mathbf{F}\) and hence to \(f(x) = \bar{f}(\phi(x)) = 0\) . By hypothesis, this implies \(x \in \mathrm{T}(\mathbf{E})\) and hence \(\phi(x) = 0\) , which proves the corollary.

COROLLARY 3. Let E be an A-module and g an A-linear mapping of E into a vector space F over K such that \(g(\mathbf{E})\) generates F and \(\operatorname{Ker}(g) \subset \mathrm{T}(\mathbf{E})\) . Then the K-linear mapping \(\bar{g}\) associated with g is an isomorphism of \(\mathrm{E}_{(\mathbb{K})}\) onto F.

\(\bar{g}\) is injective by Corollary 2 and the hypothesis that \(g(\mathbf{E})\) generates F implies that \(\bar{g}\) is surjective.

For every A-module E, the vector space \(E_{(K)}\) is said to be associated with E. For every subset S of E the rank of S over K (or by an abuse of language, the rank of S) is the rank of the canonical image \(\phi(S)\) of S in \(E_{(K)}\) , in other words (no. 2, Definition 1) the dimension over K of the vector subspace of \(E_{(K)}\) generated by \(\phi(S)\) .

When E is a torsion-free A-module, it is usually identified with its canonical image \(\phi(\mathbf{E})\) in \(\mathrm{E}_{(\mathbb{K})}\) . With this convention, every generating system of E contains a basis of \(\mathrm{E}_{(\mathbb{K})}\) (no. 1, Theorem 2). In particular:

COROLLARY 4. Every finitely generated A-module is of finite rank.

Note that the converse of this corollary is not necessarily true; for example Q is a Z-module of rank 1 but is not finitely generated over Z.

Recall (§ 5, no. 1) that for every linear mapping \(f: \mathbf{E} \to \mathbf{E}'\) (where \(\mathbf{E}\) and \(\mathbf{E}'\) are A-modules), \(f_{(\mathbf{K})}\) denotes the K-linear mapping \(1_{\mathbf{K}} \otimes f: \mathbf{E}_{(\mathbf{K})} \to \mathbf{E}_{(\mathbf{K})}'\) .

PROPOSITION 27. For every exact sequence

\[
\mathrm{E} ^ {\prime} \xrightarrow {f} \mathrm{E} \xrightarrow {g} \mathrm{E} ^ {\prime \prime}
\]

of A-linear mappings, the corresponding sequence of K-linear mappings

\[
\mathbf {E} _ {(\mathbf {K})} ^ {\prime} \xrightarrow {f _ {(\mathbf {K})}} \mathbf {E} _ {(\mathbf {K})} \xrightarrow {g _ {(\mathbf {K})}} \mathbf {E} _ {(\mathbf {K})} ^ {\prime \prime}
\]

is exact.

Suppose that \(g_{(\mathbb{K})}(\lambda^{-1} \otimes x) = 0\) , with \(\lambda \in \mathbf{A}\) , \(\lambda \neq 0\) , \(x \in \mathbb{E}\) ; this is equivalent to \(\lambda^{-1} \otimes g(x) = 0\) in \(\mathrm{E}_{(\mathbb{K})}'\) and hence also to

\[
1 \otimes g (x) = \lambda (\lambda^ {- 1} \otimes g (x)) = 0;
\]

by Proposition 26, there exists \(\alpha \neq 0\) in A such that \(\alpha g(x) = 0\) in \(\mathbf{E}^{\prime \prime}\) , or also \(g(\alpha x) = 0\) . By hypothesis, there is therefore an \(x' \in \mathbf{E}'\) such that \(\alpha x = f(x')\) and therefore \(\lambda^{-1} \otimes x = f_{(\mathbf{K})}(\alpha^{-1}\lambda^{-1} \otimes x')\) , which proves the proposition.

COROLLARY 1. If \(\mathbf{E}'\) is a submodule of \(\mathbf{E}\) , \(\mathbf{E}_{(\mathbf{K})}'\) is canonically identified with a vector subspace of \(\mathbf{E}_{(\mathbf{K})}\) and \((\mathbf{E} / \mathbf{E}')_{(\mathbf{K})}\) with \(\mathbf{E}_{(\mathbf{K})} / \mathbf{E}_{(\mathbf{K})}'\) .

It suffices to apply Proposition 27 to the exact sequence

\[
0 \rightarrow \mathrm{E} ^ {\prime} \rightarrow \mathrm{E} \rightarrow \mathrm{E} / \mathrm{E} ^ {\prime} \rightarrow 0.
\]

COROLLARY 2. For every A-linear mapping \(f: \mathrm{E} \to \mathrm{F}\) , \(\operatorname{Ker}(f_{(\mathbf{K})}) = (\operatorname{Ker}(f))_{(\mathbf{K})}\) , \(\operatorname{Im}(f_{(\mathbf{K})}) = (\operatorname{Im}(f))_{(\mathbf{K})}\) , \(\operatorname{Coker}(f_{(\mathbf{K})}) = (\operatorname{Coker}(f))_{(\mathbf{K})}\) to within canonical isomorphisms. In particular, for \(f_{(\mathbf{K})}\) to be injective (resp. surjective, resp. zero), it is necessary and sufficient that \(\operatorname{Ker}(f) \subset \mathrm{T}(\mathrm{E})\) (resp. that \(\operatorname{Coker}(f)\) be a torsion module, resp. that \(\operatorname{Im}(f) \subset \mathrm{T}(\mathrm{F})\) ).

This follows from Corollary 1 and Proposition 26.

COROLLARY 3. Let E be an A-module and \((x_{\lambda})_{\lambda \in L}\) a family of elements of E. For \((x_{\lambda})\) to be a free family, it is necessary and sufficient that in the vector K-space \(\mathbf{E}_{(\mathbb{K})}\) the family \((1 \otimes x_{\lambda})\) be free.

The family \((x_{\lambda})\) defines an A-linear mapping \(f: \mathbf{A}^{(L)} \to \mathbf{E}\) such that \(f(e_{\lambda}) = x_{\lambda}\) for all \(\lambda \in \mathbf{L}((e_{\lambda})\) being the canonical basis of \(\mathbf{A}^{(L)})\) and to say that \((x_{\lambda})\) is free means that \(f\) is injective. It suffices to apply Corollary 2 to \(f\) , observing that \(\mathbf{A}^{(L)}\) is torsion-free (Proposition 25).

\section{RESTRICTION OF THE FIELD OF SCALARS IN VECTOR SPACES}

Throughout this paragraph, K denotes a field and \(K'\) a subfield of K. On a set V, a right (resp. left) vector space structure over K defines, by restriction of scalars, a right (resp. left) vector space structure over \(K'\) .

\subsection{DEFINITION OF K'-STRUCTURES}

PROPOSITION 1. Let V be a right vector space over K and V' a subset of V which is a vector subspace over K'. The following conditions are equivalent:

\begin{enumerate}
\def\labelenumi{(\alph{enumi})}
\item
  The K-linear mapping \(\lambda\) of \(\mathrm{V}_{(\mathbb{K})}^{\prime} = \mathrm{V}^{\prime}\otimes_{\mathbb{K}^{\prime}}\mathbf{K}\) into V such that \(\lambda (x^{\prime}\otimes \xi) = x^{\prime}\xi\) for all \(x^{\prime}\in \mathrm{V}^{\prime},\xi \in \mathbf{K}\) , is bijective.
\item
  Every K'-linear mapping \(f'\) of V' into a vector K-space W can be extended uniquely to a K-linear mapping f of V into W.
\item
  Every basis of \(\mathbf{V}'\) over \(\mathbf{K}'\) is a basis of \(\mathbf{V}\) over \(\mathbf{K}\) .
\item
  There exists a basis of \(\mathbf{V}'\) over \(\mathbf{K}'\) which is also a basis of \(\mathbf{V}\) over \(\mathbf{K}\) .
\item
  The vector K-space V is generated by V' and every subset of V' free over K' is free over K.
\end{enumerate}

We know (§ 5, no. 1) that \(\mathbf{V}_{(\mathbf{K})}^{\prime}\) has a right vector K-space structure under which \((x^{\prime} \otimes \xi)\eta = x^{\prime} \otimes (\xi\eta)\) ( \(\xi, \eta\) in K, \(x^{\prime} \in \mathbf{V}^{\prime}\) ) and that for every K'-linear mapping \(f^{\prime}\) of V' into a vector K-space W, there exists one and only one K-linear mapping \(\bar{f}^{\prime}\) of \(\mathbf{V}_{(\mathbf{K})}^{\prime}\) into W such that \(\bar{f}^{\prime}(x^{\prime} \otimes 1) = f^{\prime}(x^{\prime})\) for \(x^{\prime} \in \mathbf{V}^{\prime}\) . If \(j\) is the canonical injection of V' into V, \(\lambda\) is just the corresponding K-linear mapping \(\bar{j}\) . If \(\lambda\) is bijective, then for every K'-linear mapping \(f^{\prime}: \mathbf{V}^{\prime} \to \mathbf{W}, \bar{f}^{\prime} \circ \lambda^{-1}\) is the unique K-linear mapping of V into W extending \(f^{\prime}\) ; in other words, (a) implies (b). Conversely, if (b) holds, there exists in particular a K-linear mapping \(\mu\) of V into \(\mathbf{V}_{(\mathbf{K})}^{\prime}\) such that \(\mu(x^{\prime}) = x^{\prime} \otimes 1\) for all \(x^{\prime} \in \mathbf{V}^{\prime}\) ; it is immediate that \(\mu \circ \lambda = 1_{\mathbf{V}_{(\mathbf{K})}^{\prime}}\) ; on the other hand, \(\lambda(\mu(x^{\prime})) = x^{\prime}\) for all \(x^{\prime} \in \mathbf{V}^{\prime}\) and, as by hypothesis \(j: \mathbf{V}^{\prime} \to \mathbf{V}\) can be extended uniquely to an endomorphism of V, of necessity \(\lambda \circ \mu = 1_{\mathbf{V}}\) , which completes the proof that (a) and (b) are equivalent.

For every basis \(\mathbf{B}'\) of \(\mathbf{V}'\) over \(\mathbf{K}'\) , the set \(\mathbf{B}\) of elements \(v' \otimes 1\) of \(\mathbf{V}_{(\mathbf{K})}'\) , where \(v'\) runs through \(\mathbf{B}'\) , is a basis of \(\mathbf{V}_{(\mathbf{K})}'\) over \(\mathbf{K}\) ( \(\S 5\) , no. 1, Proposition 4) and \(\lambda(\mathbf{B}) = \mathbf{B}'\) . For \(\lambda\) to be bijective, it is necessary that the image under \(\lambda\) of every basis of \(\mathbf{V}_{(\mathbf{K})}'\) be a basis of \(\mathbf{V}\) over \(\mathbf{K}\) and it is sufficient that it be so for a single basis of \(\mathbf{V}_{(\mathbf{K})}'\) ( \(\S 1\) , no. 11, Corollary 2 to Proposition 17). This proves the equivalence of (a), (c) and (d).

As every subset of \(\mathbf{V}'\) which is free over \(\mathbf{K}'\) is contained in a basis of \(\mathbf{V}'\) over \(\mathbf{K}'\) (§ 7, no. 1, Theorem 2), (c) implies (e). Finally, suppose that (e) holds; if \(\mathbf{B}'\) is a basis of \(\mathbf{V}'\) over \(\mathbf{K}'\) , it is a free subset of \(\mathbf{V}\) over \(\mathbf{K}\) ; on the other hand \(\mathbf{B}'\) generates \(\mathbf{V}'\) over \(\mathbf{K}'\) and hence generates \(\mathbf{V}\) over \(\mathbf{K}\) by hypothesis; therefore \(\mathbf{B}'\) is a basis of \(\mathbf{V}\) over \(\mathbf{K}\) , which proves that (e) implies (c).

DEFINITION 1. Let V be a right vector space over a field K and K' a subfield of K. Every vector sub-K'-space V' of V satisfying the equivalent conditions of Proposition 1 is called a K'-structure on V.

Example. Let B be a basis of V over K. For every subfield \(K'\) of K, the vector sub- \(K'\) -space of V generated by B admits B as a basis over \(K'\) and hence is a \(K'\) -structure on V. *For example, if K is commutative and V is taken to be the polynomial K-algebra \(K[X_{1},\ldots,X_{n}]\) , then, for every subfield \(K'\) of K, \(K'[X_{1},\ldots,X_{n}]\) is a \(K'\) -structure on \(V\).*

\subsection{RATIONALITY FOR A SUBSPACE}

DEFINITION 2. Let V be a right vector space over K, with a K'-structure V'. A vector of V is said to be rational over K' if it belongs to V'. A vector sub-K-space W of V is said to be rational over K' if it is generated (over K) by vectors rational over K'.

Let \((v_{\iota}^{\prime})_{\iota\in I}\) be a basis of \(V^{\prime}\) over \(K^{\prime}\) , which is therefore also a basis of \(V\) over \(K\) (no. 1, Proposition 1). For a vector \(x = \sum_{\iota} v_{\iota}^{\prime} \xi_{\iota}\) of \(V\) to be rational over \(K^{\prime}\) , it is necessary and sufficient that \(\xi_{\iota} \in K^{\prime}\) for all \(\iota \in I\) .

If W is a vector sub-K-space of V which is rational over \(K'\) , it follows from Definition 2 that \(W' = W \cap V'\) is a vector sub- \(K'\) -space of W, which generates W over K; on the other hand every subset of \(W'\) which is free over \(K'\) is also free over K since it is contained in \(V'\) (no. 1, Proposition 1). It follows (no. 1, Proposition 1) that \(W'\) is a \(K'\) -structure on W said to be induced by the \(K'\) -structure \(V'\) on V.

For every vector sub-K'-space W' of V', we shall denote by W'.K the vector sub-K-space of V consisting of the linear combinations of elements of W' with coefficients in K.

PROPOSITION 2. Let V be a right vector space over K and V' a K'-structure on V. The mapping W' → W'.K is a bijection of the set of vector sub-K'-spaces of V' onto the set of vector sub-K-spaces of V which are rational over K' and the inverse bijection is W → W ∩ V'.

Clearly the bijection \(\lambda^{-1}: V \to V' \otimes_{K'} K\) , the inverse of the bijection \(\lambda\) defined in no. 1, Proposition 1, maps every vector sub- \(K'\) -space \(W'\) of \(V'\) onto its image under the canonical injection \(x' \mapsto x' \otimes 1\) and \(W'. K\) onto \(W' \otimes_{K'} K\) ; the assertions of Proposition 2 are thus consequences of Definition 2 and §7, no. 9, Proposition 19.

COROLLARY 1. Every sum and every intersection of vector sub-K-spaces of V which are rational over K' is a rational subspace over K'.

The assertion concerning sums is obvious. On the other hand, if \((\mathbf{W}_{\iota}')_{\iota \in I}\) is a family of vector sub- \(\mathbf{K}'\) -spaces of \(\mathbf{V}'\) , then \(\left(\bigcap_{\iota \in I} \mathbf{W}_{\iota}'\right) \otimes_{\mathbb{K}'} \mathbf{K} = \bigcap_{\iota \in I} (\mathbf{W}_{\iota}' \otimes_{\mathbb{K}'} \mathbf{K})\) (§ 7, no. 7, Corollary to Proposition 14), which proves the corollary.

A basis B of V over K is said to be rational over \(K'\) if it consists of vectors which are rational over \(K'\) .

COROLLARY 2. Every basis of V over K which is rational over \(\mathbf{K}'\) is a basis of \(\mathbf{V}'\) over \(\mathbf{K}'\) .

If \(W'\) is the vector sub- \(K'\) -space of \(V'\) generated by B, then \(W'. K = V = V'. K\), whence \(V' = W'\) by virtue of Proposition 2.

\subsection{RATIONALITY FOR A LINEAR MAPPING}

DEFINITION 3. Let \(\mathbf{V}_1, \mathbf{V}_2\) be two right vector spaces over \(\mathbf{K}\) with \(\mathbf{K}'\) -structures \(\mathbf{V}_1'\) , \(\mathbf{V}_2'\) respectively. A K-linear mapping \(f: \mathbf{V}_1 \to \mathbf{V}_2\) is said to be rational over \(\mathbf{K}'\) if \(f(\mathbf{V}_1') \subset \mathbf{V}_2'\) .

If \(V_{3}\) is a third right vector space over K, with a \(K^{\prime}\) -structure \(V_{3}^{\prime}\) and a K-linear mapping \(g:V_{2}\rightarrow V_{3}\) is rational over \(K^{\prime}\) , clearly \(g\circ f:V_{1}\rightarrow V_{3}\) is rational over \(K^{\prime}\) .

PROPOSITION 3. Let \(\mathbf{V}_1, \mathbf{V}_2\) be two right vector spaces over \(\mathbf{K}\) and \(\mathbf{V}_1', \mathbf{V}_2' \mathbf{K}'\) -structures on \(\mathbf{V}_1, \mathbf{V}_2\) respectively. \(\mathbf{V}_1\) (resp. \(\mathbf{V}_2\) ) is canonically identified with \(\mathbf{V}_1' \otimes_{\mathbb{K}'} \mathbf{K}\) (resp. \(\mathbf{V}_2' \otimes_{\mathbb{K}'} \mathbf{K}\) ) (no. 1, Proposition 1).

\begin{enumerate}
\def\labelenumi{(\roman{enumi})}
\item
  The mapping \(f' \mapsto f' \otimes 1_{\mathbf{K}} = f_{(\mathbf{K})}'\) is a bijection of \(\operatorname{Hom}_{\mathbf{K}'}(\mathbf{V}_1', \mathbf{V}_2')\) onto the set of K-linear mappings of \(\mathbf{V}_1\) into \(\mathbf{V}_2\) which are rational over \(\mathbf{K}'\) ; the inverse bijection associates with every K-linear mapping \(f: \mathbf{V}_1 \to \mathbf{V}_2\) rational over \(\mathbf{K}'\) the K'-linear mapping \(f': \mathbf{V}_1' \to \mathbf{V}_2'\) with the same graph as the restriction off to \(\mathbf{V}_1'\) .
\item
  For every K-linear mapping \(f: \mathbf{V}_1 \to \mathbf{V}_2\) rational over \(\mathbf{K}'\) ,
\end{enumerate}

\[
f (\mathrm{V} _ {1} ^ {\prime}) = f (\mathrm{V} _ {1}) \cap \mathrm{V} _ {2} ^ {\prime} \quad and \quad \bar {f} ^ {- 1} (\mathrm{V} _ {2} ^ {\prime}) = \mathrm{V} _ {1} ^ {\prime} + \operatorname{Ker} (f).
\]

\begin{enumerate}
\def\labelenumi{(\roman{enumi})}
\item
  Clearly with the above identifications, if \(f': \mathbf{V}_1' \to \mathbf{V}_2'\) is a K'-linear mapping, \(f_{(\mathbf{K})}' = f' \otimes 1_{\mathbf{K}}\) is rational over \(\mathbf{K}'\) and \(f'\) is the mapping with the same graph as the restriction of \(f_{(\mathbf{K})}'\) to \(\mathbf{V}_1'\) . Conversely, if \(f: \mathbf{V}_1 \to \mathbf{V}_2\) is a K-linear mapping rational over \(\mathbf{K}'\) and \(f': \mathbf{V}_1' \to \mathbf{V}_2'\) has the same graph as the restriction of \(f\) to \(\mathbf{V}_1'\) , \(f\) and \(f_{(\mathbf{K})}'\) coincide on \(\mathbf{V}_1'\) , which is a generating system of \(\mathbf{V}_1\) over \(\mathbf{K}\) , hence \(f = f_{(\mathbf{K})}'\) .
\item
  If \(f = f' \otimes 1_{\mathbb{K}}\) , then \(f(\mathbf{V}_1) = f(\mathbf{V}_1' \otimes_{\mathbb{K}'} \mathbf{K}) = f'(\mathbf{V}_1') \otimes_{\mathbb{K}'} \mathbf{K}\) and, as \(f'(\mathbf{V}_1') \subset \mathbf{V}_2'\) , \(f(\mathbf{V}_1') = f'(\mathbf{V}_1') = f(\mathbf{V}_1) \cap \mathbf{V}_2'\) (§ 7, no. 9, Proposition 19); the formula \(f^{-1}(\mathbf{V}_2') = \mathbf{V}_1' + \operatorname{Ker}(f)\) follows immediately.
\end{enumerate}

COROLLARY 1. In the notation of Proposition 3, \(\operatorname{Im}(f) = (\operatorname{Im}(f'))_{(\mathbb{K})}\) ,

\[
\operatorname{Ker} (f) = (\operatorname{Ker} (f ^ {\prime})) _ {(\mathbb {K})}, \quad \operatorname{Coker} (f) = (\operatorname{Coker} (f ^ {\prime})) _ {(\mathbb {K})}.
\]

In particular, for \(f\) to be injective (resp. surjective, zero), it is necessary and sufficient that \(f'\) be so. If \(f\) is bijective, its inverse mapping is rational over \(\mathbf{K}'\) .

This is a particular case of § 7, no. 9, Corollary to Proposition 18.

COROLLARY 2. Let \(f: \mathbf{V}_1 \to \mathbf{V}_2\) be a K-linear mapping rational over \(\mathbf{K}'\) . For every vector sub-K-space \(\mathbf{W}_1\) of \(\mathbf{V}_1\) (resp. \(\mathbf{W}_2\) of \(\mathbf{V}_2\) ) rational over \(\mathbf{K}'\) , \(f(\mathbf{W}_1)\) (resp. \(f(\mathbf{W}_2)\) ) is a vector sub-K-space of \(\mathbf{V}_2\) (resp. \(\mathbf{V}_1\) ) rational over \(\mathbf{K}'\) .

In the notation of Proposition 3, for every vector sub-K'-space W \(_{1}\) of V \(_{1}\) , f \(_{(K)}\) (W \(_{1}\) ' ⊗ \(_{K'}\) K) = f'(W \(_{1}\) ) ⊗ \(_{K'}\) K; whence the assertion relating to W \(_{1}\) (no. 2, Proposition 2). On the other hand, let W \(_{2}\) be a vector sub-K'-space of V \(_{2}\) and let g' be the canonical K'-linear mapping V \(_{2}\) → V \(_{2}\) /W \(_{2}\) ; then \(f'^{-1}(W_{2}') = \text{Ker}(g' \circ f')\) ; hence, by Corollary 1,

\[
\bar {f} _ {(\mathbf {K})} ^ {\prime - 1} (\mathbf {W} _ {2} ^ {\prime} \otimes_ {\mathbf {K} ^ {\prime}} \mathbf {K}) = \bar {f} ^ {\prime - 1} (\mathbf {W} _ {2} ^ {\prime}) \otimes_ {\mathbf {K} ^ {\prime}} \mathbf {K},
\]

whence the assertion relating to \(W_{2}\) .

Let \(V_{1}, V_{2}\) be two right vector K-spaces with \(K'\) -structures \(V_{1}', V_{2}'\) respectively. It is immediate that \(V_{1}' \times V_{2}'\) is a \(K'\) -structure on \(V_{1} \times V_{2}\) , called the product of the \(K'\) -structures \(V_{1}'\) and \(V_{2}'\) .

PROPOSITION 4. For a K-linear mapping \(f: \mathbf{V}_1 \to \mathbf{V}_2\) to be rational over \(\mathbf{K}'\) , it is necessary and sufficient that its graph \(\Gamma\) be rational over \(\mathbf{K}'\) for the product \(\mathbf{K}'\) -structure on \(\mathbf{V}_1 \times \mathbf{V}_2\) .

Let \(g\) be the mapping \(x_1 \mapsto (x_1, f(x_1))\) from \(V_1\) to \(V_1 \times V_2\) ; this is a K-linear mapping such that \(\Gamma = g(V_1)\) ; if \(f\) is rational over \(K'\) , so is \(\Gamma\) by virtue of Corollary 2 to Proposition 3. Conversely, suppose that \(\Gamma\) is rational over \(K'\) and let it have the \(K'\) -structure induced by that on \(V_1 \times V_2\) ; it follows immediately from the definitions that the restrictions \(p_1, p_2\) to \(\Gamma\) of the projections \(pr_1, pr_2\) , are K-linear mappings, rational over \(K'\) of \(\Gamma\) into \(V_1\) and \(V_2\) respectively. As \(p_1\) is bijective, its inverse mapping \(q_1\) is rational over \(K'\) (Corollary 1 to Proposition 3) and hence so is \(f = p_2 \circ q_1\) .

\subsection{RATIONAL LINEAR FORMS}

Let V be a right vector space over K with a \(K^{\prime}\) -structure \(V^{\prime}\) . As \(K_{d}^{\prime}\) is a \(K^{\prime}\) -structure on the right vector K-space \(K_{d}\) , we may define linear forms \(x^{*} \in V^{*}\) , rational over \(K^{\prime}\) , as the linear mappings of V into \(K_{d}\) , rational over \(K^{\prime}\) for the

K'-structures on V and \(K_{d}\) . By virtue of no. 3, Proposition 3, the set \(R'\) of these linear forms is the image of the dual \(V'^*\) of \(V'\) under the composite mapping

\[
\mathrm{V} ^ {\prime *} \xrightarrow {\phi} \mathrm{K} \otimes_ {\mathrm{K} ^ {\prime}} \mathrm{V} ^ {\prime *} \xrightarrow {\upsilon} \mathrm{V} ^ {*}\tag{1}
\]

where \(\phi(x^{\prime*}) = 1 \otimes x^{\prime*}\) and \(\upsilon(\xi \otimes x^{\prime*})\) is the linear form \(y^{*}\) on V such that \(y^{*}(x^{\prime}) = \xi\langle x^{\prime*}, x^{\prime} \rangle\) for all \(x^{\prime} \in V^{\prime}\) (§ 5, no. 4). We know that this mapping is injective (§ 7, no. 9, Propositions 20 and 19) and clearly \(R^{\prime}\) is a left vector sub- \(K^{\prime}\) -space of \(V^{*}\) ; moreover every subset of \(R^{\prime}\) free over \(K^{\prime}\) is free over K. But in general \(R^{\prime}\) does not necessarily generate \(V^{*}\) over K and does not therefore define a \(K^{\prime}\) -structure on \(V^{*}\) (Exercise 2). However, if V is of finite dimension \(n\) over K, \(V^{\prime*}\) is of dimension \(n\) over \(K^{\prime}\) and \(R^{\prime}\) then defines canonically a \(K^{\prime}\) -structure on \(V^{*}\) .

PROPOSITION 5. Let V be a right vector space over K, V' a K'-structure on V and W a vector sub-K-space of V. For W to be rational over K', it is necessary and sufficient that there exist a set H ⊂ V* of rational linear forms over K' such that W is the orthogonal of H in V (§ 2, no. 4).

Let H be a subset of V* whose elements are rational linear forms over K'. For all \(x^{*} \in H\) , the kernel of \(x^{*}\) is a vector sub-K-space of V, rational over K' (no. 3, Corollary 2 to Proposition 3); the intersection of these kernels is therefore also a vector sub-K-space of V, rational over K' (no. 2, Corollary 1 to Proposition 2).

Conversely, let W be a vector sub-K-space of V rational over \(K'\) , so that W is identified with \(W' \otimes_{K'} K\) , where \(W' = W \cap V'\) (no. 2, Proposition 2). For a linear form \(x'^* \in V'^*\) to be zero on \(W'\) , it is necessary and sufficient that the linear form \(x^* \in V^*\) which corresponds to it under (1) be zero on W, for by no. 3, Corollary 1 to Proposition 3,

\[
\operatorname{Ker} (x ^ {*}) = \left(\operatorname{Ker} (x ^ {\prime *})\right) \otimes_ {\mathbb {K} ^ {\prime}} \mathrm{K} \quad \text { and } \quad \operatorname{Ker} (x ^ {* ^ {\prime}}) = (\operatorname{Ker} (x ^ {*})) \cap \mathrm{V} ^ {\prime}.
\]

Let \(H'\) be the orthogonal of \(W'\) in \(V'^*\); we know (§ 7, no. 5, Theorem 7) that \(W'\) is the orthogonal of \(H'\) in \(V'\) ; if \(H\) is the image of \(H'\) in \(V*\) under the mapping (1), it follows from the above that \(W\) is the orthogonal of \(H\) in \(V\) , taking account of no. 7, Corollary to Proposition 14.

\subsection{APPLICATION TO LINEAR SYSTEMS}

PROPOSITION 6. (i) Given a system of homogeneous linear equations

\[
\sum_ {\iota \in I} \alpha_ {\mu \iota} \xi_ {\iota} = 0 (\mu \in M)\tag{2}
\]

whose coefficients \(\alpha_{\mu \iota}\) belong to \(K'\) , every solution ( \(\xi_{\iota}\) ) of this system consisting of elements of K is a linear combination with coefficients in K of solutions ( \(\xi_{\iota}'\) ) of (2) consisting of elements of \(K'\) .

\begin{enumerate}
\def\labelenumi{(\roman{enumi})}
\setcounter{enumi}{1}
\tightlist
\item
  Given a system of linear equations
\end{enumerate}

\[
\sum_ {\iota \in I} \alpha_ {\mu \iota} \xi_ {\iota} = \beta_ {\mu} \quad (\mu \in \mathbf {M})\tag{3}
\]

whose coefficients \(\alpha_{\mu \iota}\) and right hand sides \(\beta_{\mu}\) belong to \(K'\) , if there exists a solution to the system consisting of elements of K, there also exists a solution consisting of elements of \(K'\) .

\begin{enumerate}
\def\labelenumi{(\roman{enumi})}
\item
  For every set S, let the right vector K-space \(\mathrm{K}_{d}^{(\mathrm{S})}\) be given the \(K^{\prime}\) -structure \(\mathrm{K}_{d}^{\prime(\mathrm{S})}\) . Let f be the K-linear mapping of \(\mathrm{K}_{d}^{(1)}\) into \(\mathrm{K}_{d}^{(\mathrm{M})}\) mapping every vector \((\xi_{\iota})_{\iota\in\mathbb{I}}\) to the vector \((\zeta_{\mu})_{\mu\in\mathbb{M}}\) defined by \(\zeta_{\mu}=\sum_{\iota\in\mathbb{I}}\alpha_{\mu\iota}\xi_{\iota}\) for all \(\mu\in M\) . Clearly f is rational over \(K^{\prime}\) ; its kernel V, which is the set of solutions in K of the system (2), is a subspace of \(\mathrm{K}_{d}^{(\mathrm{I})}\) rational over \(K^{\prime}\) (no. 3, Corollary 2 to Proposition 3) and hence generated by the solutions of (2) in \(K^{\prime}\) .
\item
  We consider K as a left vector \(\mathbf{K}'\) -space; there exists a \(\mathbf{K}'\) -linear projector \(p\) of K onto its vector subspace \(\mathbf{K}_s'\) (§ 7, no. 3, Proposition 4); if \((\xi_{\iota})\) is a solution of (3) in K, then \(\sum_{\iota \in I} \alpha_{\mu \iota} p(\xi_{\iota}) = p\left(\sum_{\iota \in I} \alpha_{\mu \iota} \xi_{\iota}\right) = p(\beta_\mu) = \beta_\mu\) , which proves that \((p(\xi_{\iota}))\) is a solution of (3) in \(\mathbf{K}'\) .
\end{enumerate}

A ring K is called (left) faithfully flat over a subring \(K'\) if Proposition 6 holds for K and \(K'\) ; we shall study this notion in detail later (Commutative Algebra, I, §3).

\subsection{SMALLEST FIELD OF RATIONALITY}

Let V be a right vector K-space with a K'-structure V'. For every field L such that \(K' \subset L \subset K\) , we write \(V_{L} = V'. L\); clearly every basis of \(V'\) over \(K'\) is a basis of V over K and a basis of \(V_{L}\) over L. Hence \(V_{L}\) is an L-structure on V and \(V'\) a K'-structure on \(V_{L}\) .

PROPOSITION 7. (i) Let V be a right vector K-space with a K'-structure V'. For every vector \(x \in V\) (resp. every vector sub-K-space W of V), the set of subfields L of K containing K' and such that x (resp. W) is rational over L has a least element \(\mathrm{K}'(x)\) (resp. \(\mathrm{K}'(\mathrm{W})\) ).

\begin{enumerate}
\def\labelenumi{(\roman{enumi})}
\setcounter{enumi}{1}
\tightlist
\item
  Let \(V_{1}\) , \(V_{2}\) be two right vector K-spaces with \(K^{\prime}\) -structures \(V_{1}^{\prime}\) , \(V_{2}^{\prime}\) respectively. For every K-linear mapping f of \(V_{1}\) into \(V_{2}\) , the set of subfields L of K containing \(K^{\prime}\) and such that f is rational over L has a least element \(\mathrm{K}^{\prime}(f)\) .
\end{enumerate}

We first prove assertion (i) for a vector \(x \in V\) . Let \(B\) be a basis of \(V\) rational over \(K'\) ; \(B\) is a basis of \(V'\) over \(K'\) and a basis of \(V_L\) over \(L\) for every field \(L\) such that \(K' \subset L \subset K\) ; for \(x = \sum_{b \in B} b\xi_b\) to be rational over \(L\) , it is necessary and sufficient that the \(\xi_{b}\) belong to L (no. 2) and hence the smallest field L with this property is the field generated by K', and the \(\xi_{b}\) for \(b \in \mathbf{B}\) .

We next prove (ii). Let \(\mathbf{B}_1, \mathbf{B}_2\) be bases of \(\mathbf{V}_1, \mathbf{V}_2\) respectively, which are rational over \(\mathbf{K}'\) and write, for all \(b_1 \in \mathbf{B}_1\) , \(f(b_1) = \sum_{b_2 \in \mathbf{B}_2} b_2 \alpha_{b_2 b_1}\) (*the family \((\alpha_{b_2 b_1})\) is just the matrix of \(f\) with respect to the bases \(\mathbf{B}_1\) and \(\mathbf{B}_2\) ; cf.~§ 10, no. 4*). As \(\mathbf{B}_1\) (resp. \(\mathbf{B}_2\) ) is a basis of \((\mathbf{V}_1)_L\) (resp. \((\mathbf{V}_2)_L\) ) over \(L\) for every field such that \(\mathbf{K}' \subset L \subset K\) , for \(f\) to be rational over \(L\) , it is necessary and sufficient that the \(\alpha_{b_2 b_1}\) belong to \(L\) ; the smallest field with this property is therefore the field generated by \(\mathbf{K}'\) and the \(\alpha_{b_2 b_1}\) for \(b_1 \in \mathbf{B}_1\) , \(b_2 \in \mathbf{B}_2\) .

Finally, to establish assertion (i) for a subspace W of V, we shall first prove the following lemma:

Lemma 1. Let V be a right vector K-space with a K'-structure V' and W a vector sub-K-space of V. There exist two vector sub-K-spaces W₁, W₂ of V, rational over K', such that V is the direct sum of W₁ and W₂ and, if V is identified with W₁ × W₂, W is the graph of a K-linear mapping g of W₁ into W₂.

Let B be a basis of V rational over \(K'\) . Applying Theorem 2 of §7, no. 1, to a basis of W over K, considered as a free subset of V, and the generating system the union of this free subset and B, it is seen that there exists a subset C of B such that V is the direct sum of W and the subspace \(W_{2}\) of V generated by C. Also let \(W_{1}\) be the subspace of V generated by B - C. As \(B \subset V'\) , clearly \(W_{1}\) and \(W_{2}\) are rational over \(K'\) . Moreover, for all \(x \in W_{1}\) , there exists one and only one vector \(g(x)\) of \(W_{2}\) such that \(x + g(x) \in W\) , since V is the direct sum of W and \(W_{2}\) ; then W is the graph of g and g is K-linear since W is a vector sub-K-space of V.

Having proved this lemma, it is known that W is rational over a subfield L of K containing \(K'\) if and only if g is rational over L (no. 3, Proposition 4). The smallest field \(\mathrm{K}'(g)\) such that g is rational over \(\mathrm{K}'(g)\) is therefore also the smallest field over which W is rational.
